%% 4-10 向きと大きさをもつ量：ベクトル（第13回：組 V4-6。ページ制限なし。偶数ページ）
%% ほかの巻・物理シリーズから指される形：
%%   ・ベクトルの和（矢印をつなぐ／平行四辺形）、合力（2-3 の表「ベクトルの合成（4-10）」）
%%   ・-a は反対の向き（2-2）、平行移動（3-2 の「右へ3」）
%%   ・差。速さが同じでも向きが変われば Δv≠0（物理の円運動）、相対速度＝相手の速度－自分の速度
%% 数（合力・Δv・相対速度・クレーンのつり角）は、すべて Python で計算して確かめた（scratchpad/v4-6/py/c10.py）。
\makeatletter
% 説明の段の見出し（丸数字＋太字の1行）  \bkivkstep{番号}{見出し}
\newcommand\bkivkstep[2]{\par\addvspace{6pt}\Needspace{5\baselineskip}\noindent
  {\bkcircled[bkOrange]{#1}\enspace\sffamily\bfseries\color{bkOrange!75!black}#2}\par\nobreak\vspace{2pt}}
% 式を、まん中に1行で置く  \bkivkline{…}
\newcommand\bkivkline[1]{\par\addvspace{3pt}{\centering #1\par}\addvspace{3pt}\noindent\ignorespaces}
% 小見出し（それだけで1行）  \bkivkmidashi{…}
\newcommand\bkivkmidashi[1]{\par\addvspace{4pt}\Needspace{5\baselineskip}\noindent\textbf{#1}\par\nobreak}
% 荷物（上から見た段ボール箱）  \bkivkbox{(中心)}{1辺 mm}
\newcommand\bkivkbox[2]{\begin{scope}[shift={#1}]
  \filldraw[fill=bkOrange!28, draw=bkOrange!70!black, line width=0.7pt, rounded corners=0.6mm] (-#2/2,-#2/2) rectangle (#2/2,#2/2);
  \draw[bkOrange!60!black, line width=0.5pt] (-#2/2,0) -- (#2/2,0);\end{scope}}
% 方位（北と東）  \bkivkcompass{(x,y)}
\newcommand\bkivkcompass[1]{\begin{scope}[shift={#1}]
  \draw[bkGray, line width=0.7pt, -{Stealth[length=3pt]}] (0,0) -- (0,7);
  \draw[bkGray, line width=0.7pt, -{Stealth[length=3pt]}] (0,0) -- (7,0);
  \node[font=\sffamily\footnotesize, anchor=south, inner sep=1pt] at (0,7.4) {北};
  \node[font=\sffamily\footnotesize, anchor=west, inner sep=1pt] at (7.4,0) {東};\end{scope}}
% ばねばかり（引く向きの角度 #2 にそって置く）  \bkivkscale{(取っ手の先)}{角度}
\newcommand\bkivkscale[2]{\begin{scope}[shift={#1}, rotate=#2]
  \filldraw[fill=bkBlue!12, draw=bkBlue!70!black, line width=0.6pt, rounded corners=0.8mm] (-13,-2.2) rectangle (-1,2.2);
  \foreach \k in {0,...,5} {\draw[bkBlue!70!black, line width=0.3pt] ({-11.5+\k*1.8},2.2) -- ({-11.5+\k*1.8},1.2);}
  \draw[bkInk, line width=0.6pt] (-1,0) -- (1.5,0);
  \draw[bkInk, line width=0.6pt] (2.6,0) circle (1.1);\end{scope}}
\makeatother
\tikzset{kivk vec/.style={-{Stealth[length=4.5pt, width=4pt]}, line width=1.5pt},
  kivk lab/.style={font=\sffamily\footnotesize, inner sep=1pt}}

\bkchapter{4}{10}{向きと大きさをもつ量：\\ベクトル}{\bklevel{2}}{向きと大きさをもつ量：ベクトル}
\bkmichishirube{2-2,3-2,4-4,4-7}{4-11,7-2,8-5}

\begin{goal}
\begin{itemize}[label={\color{\bkTT}●}, leftmargin=1.6\zw, itemsep=0pt, topsep=1pt]
\item \bklevel{2}\enspace 向きと大きさをもつ量を矢印で表し、$\vec{a}$・$\overrightarrow{\mathrm{AB}}$・$|\vec{a}|$ の記号で書ける。
\item \bklevel{2}\enspace 2つの力や移動を、矢印をつないで（平行四辺形で）たし、合わせた大きさを求められる。
\item \bklevel{2}\enspace ベクトルの実数倍と差をかき、速度の変化や相対速度を求められる。
\end{itemize}
\end{goal}

\begin{hatena}
\noindent\begin{minipage}[c]{0.54\linewidth}\raggedright 
ソラさんとミナさんが、床の上の重い荷物をロープで引きます。上から見た図のように、ソラさんは東へ\qty{40}{N}、ミナさんは北へ\qty{30}{N}の力で引きます。

荷物は、どっちへ動き出すかな？　東かな、北かな、それとも、そのあいだかな？　2人の力を合わせると、$40+30=70$ で\qty{70}{N}になるのかな？
\end{minipage}\hfill
\begin{minipage}[c]{0.43\linewidth}
\centering
\begin{tikzpicture}[x=0.85mm, y=0.85mm]
  \bkivkbox{(0,0)}{12}
  \draw[kivk vec, bkBlue] (6.5,0) -- (32,0);
  \draw[kivk vec, bkOrange] (0,6.5) -- (0,25.5);
  \node[kivk lab, text=bkBlue!80!black, anchor=north] at (19,-1.5) {ソラ　\qty{40}{N}};
  \node[kivk lab, text=bkOrange!80!black, anchor=east] at (-1.5,17) {ミナ　\qty{30}{N}};
  \bkkid{(37,-5)}{14}{bkBlue!70}{bkInk}
  \bkkid{(0,28)}{14}{bkOrange!80}{bkInk}
  \draw[bkGray, line width=0.9pt, dashed, -{Stealth[length=4pt]}] (5,3.75) -- (24,18);
  \node[font=\sffamily\bfseries\large, text=bkGray] at (27,21.5) {?};
  \bkivkcompass{(30,31)}
\end{tikzpicture}
\end{minipage}
\end{hatena}

\noindent 長さや時間は、「\qty{3}{m}」「\qty{5}{s}」のように、数と単位だけで言いつくせます。ところが、力はちがいます。同じ\qty{40}{N}でも、東へ引くか北へ引くかで、荷物の動く向きが変わります。力には、「どれだけの強さか」のほかに、「どっちを向いているか」があるのです。

この章では、向きと大きさをもつ量を\textbf{矢印}で表します。そして、矢印どうしをたしたり、ひいたりする方法を作ります。この矢印の計算は、力・速度・電場や磁場など、物理と工学のほとんどすべての場面で使われます。

%% ---------------- やってみよう・思い出そう ----------------
\begin{yattemiyou}[3つのばねばかりで、輪を引き合おう]
\textbf{用意するもの}\par
ばねばかり3つ（\qty{10}{N}まではかれるもの）、小さな金属の輪（キーホルダーの輪など）、糸、画用紙、定規、分度器
\par\vspace{3pt}
\noindent\begin{minipage}[c]{0.55\linewidth}\raggedright 
\begin{enumerate}[label={\bkcircled[bkGreen]{\arabic*}}, leftmargin=2.4\zw, itemsep=1pt, topsep=0pt]
\item 画用紙を机にテープでとめ、まん中に輪を置く。輪に糸を3本結び、それぞれの先にばねばかりをつける。
\item 3人で、ばねばかりを机と平行に引く。輪が止まったら、3つの目もりを読み、糸の向きにそって鉛筆で線を引く。
\item 輪の中心から、糸の向きに、\qty{1}{N}を\qty{1}{cm}として矢印をかく。
\item 矢印AとBを2辺とする平行四辺形をかき、対角線の長さと向きを、矢印Cとくらべよう。
\end{enumerate}
\end{minipage}\hfill
\begin{minipage}[c]{0.42\linewidth}\centering
\begin{tikzpicture}[x=0.88mm, y=0.88mm]
  \fill[bkPaper] (-27,-27) rectangle (32,30);
  \draw[bkLine, line width=0.4pt] (-27,-27) rectangle (32,30);
  % 糸（A：東、B：北、C：南西よりの向き 233.13°）
  \draw[bkInk!70, line width=0.5pt] (1.6,0) -- (13,0);
  \draw[bkInk!70, line width=0.5pt] (0,1.6) -- (0,11);
  \draw[bkInk!70, line width=0.5pt] (233.13:1.6) -- (233.13:12);
  \bkivkscale{(16.6,0)}{180}
  \bkivkscale{(0,14.6)}{270}
  \bkivkscale{($(233.13:15.6)$)}{53.13}
  \draw[bkInk, line width=0.9pt] (0,0) circle (1.6);
  \node[font=\sffamily\bfseries\small] at (22,4.5) {A};
  \node[font=\sffamily\bfseries\small] at (4.5,22) {B};
  \node[font=\sffamily\bfseries\small] at (-19,-11) {C};
\end{tikzpicture}
\end{minipage}
\par\vspace{4pt}
ソラさんたちの結果です。
\begin{bkcenter}
\begin{bktable}{colspec={Q[l,m,wd=30mm] Q[c,m,wd=24mm] Q[c,m,wd=24mm] Q[c,m,wd=24mm]}, rowsep=1.5pt, row{1}={font=\sffamily\small}}
 & A & B & C \\
目もり & \qty{3.0}{N} & \qty{4.0}{N} & \qty{5.0}{N} \\
矢印の長さ & \qty{3.0}{cm} & \qty{4.0}{cm} & \qty{5.0}{cm} \\
\end{bktable}
\bktabcaption{AとBの糸のあいだの角は90°だった。}
\end{bkcenter}
\end{yattemiyou}

\begin{omoidasou}
\bkrecall{2-2}{マイナスは「反対の向き」を表す。東を＋と決めると、西へ\qty{3}{m}は $-\qty{3}{m}$。＋の向きは、はじめに決めて変えない。}
\bkrecall{2-15}{力の大きさは N（ニュートン）ではかる。質量約\qty{100}{g}の物にはたらく重さ（重力）が約\qty{1}{N}。}
\bkrecall{3-2}{点 $(a,\ b)$ を右へ $p$、上へ $q$ 動かすと $(a+p,\ b+q)$。どれだけ動いたかは「あと－まえ」で求める。}
\bkrecall{4-4}{直角をはさむ2辺が $a$、$b$、斜辺が $c$ の直角三角形では、$a^2+b^2=c^2$（三平方の定理）。}
\bkrecall{4-6}{直角三角形で、$\tan\theta$＝向かいの辺÷となりの辺。電卓の $\tan^{-1}$ で、$\tan$ の値から角を求められる。}
\bkrecall{4-7}{余弦定理：三角形の2辺 $b$、$c$ と、そのあいだの角 $A$ から、向かいの辺 $a$ は $a^2=b^2+c^2-2bc\cos A$。鈍角の $\cos$ は負（$\cos120°=-0.5$）。}
\end{omoidasou}

\noindent やってみようの答え：AとBを2辺とする長方形（平行四辺形の1つ）の対角線は、三平方の定理（4-4）で $\sqrt{3^2+4^2}=\sqrt{25}=5$ なので\qty{5.0}{cm}です。Cの矢印と同じ長さで、向きはちょうど反対になりました。AとBの2つの力を合わせると、\qty{5.0}{N}の力1つと同じ働きをします。その働きを、Cがちょうど打ち消していたので、輪は止まっていたのです。$3+4=7$ の\qty{7}{N}ではないことに注意しましょう。

\bkfukidashi{sora}{3 N と 4 N を合わせたのに、7 N じゃなくて 5 N なんだ。向きがちがうと、たし算のしかたも変わるんだね。}

%% ---------------- 数学の道具 ----------------
\bksection{doguu}{数学の道具}

\bksubsection{1}{「どっちへ」がいる量と、いらない量}

\label{bk:4-10:vector}駅で待ち合わせをした友だちから、メッセージが来ました。「駅から\qty{300}{m}歩いた所にいるよ」。どこをさがせばいいかな？

\begin{bkfigure}
\begin{tikzpicture}[x=1mm, y=1mm]
  % 左：300 m だけ
  \draw[bkGray, line width=0.7pt, dashed] (0,0) circle (18);
  \foreach \a in {20,75,140,200,250,320} {\fill[bkOrange] (\a:18) circle (1.3);}
  \foreach \a in {20,140,250} {\node[font=\sffamily\bfseries\small, text=bkOrange!80!black] at (\a:23) {?};}
  \filldraw[fill=bkBlue!20, draw=bkBlue!80!black, line width=0.6pt] (-4,-2.5) rectangle (4,2.5);
  \node[font=\sffamily\footnotesize] at (0,0) {駅};
  \draw[bkGray, line width=0.5pt, <->] (4.5,0) -- (17.5,0);
  \node[font=\sffamily\footnotesize, anchor=south, fill=white, inner sep=0.8pt] at (11,0.8) {300 m};
  \node[font=\sffamily\small, anchor=north] at (0,-26) {「駅から\qty{300}{m}」};
  % 右：東へ 300 m
  \begin{scope}[shift={(66,0)}]
  \filldraw[fill=bkBlue!20, draw=bkBlue!80!black, line width=0.6pt] (-4,-2.5) rectangle (4,2.5);
  \node[font=\sffamily\footnotesize] at (0,0) {駅};
  \draw[kivk vec, bkRed] (4.5,0) -- (18,0);
  \fill[bkOrange] (19.5,0) circle (1.3);
  \node[font=\sffamily\footnotesize, anchor=west] at (21.5,0) {ここ！};
  \node[font=\sffamily\footnotesize, anchor=south, text=bkRed!80!black] at (11,1.5) {東へ300 m};
  \node[font=\sffamily\small, anchor=north] at (6,-26) {「駅から東へ\qty{300}{m}」};
  \bkivkcompass{(-6,12)}
  \end{scope}
\end{tikzpicture}
\bkfigcaption{「\qty{300}{m}」だけでは、駅を中心とする円の上のどこかしかわからない。向きをそえると、1つの点に決まる。}
\end{bkfigure}

「\qty{300}{m}」だけでは、駅を中心とする半径\qty{300}{m}の円の上の、どこにいるかわかりません。「東へ\qty{300}{m}」と向きもそえると、場所はただ1つに決まります。\textbf{どこからどこへ動いたか}を伝えるには、「どれだけ」と「どっちへ」の2つが必要なのです。

力も同じです。机の上の消しゴムを、同じ強さで右から押すと左へ、手前から押すとおくへ動きます。力の働きを伝えるにも、「どれだけ」と「どっちへ」の2つが必要です。

身のまわりの量を、2つに分けてみましょう。

\begin{bkcenter}
\begin{bktable}{colspec={Q[l,m,wd=40mm] Q[l,m,wd=80mm]}, rowsep=2pt}
量の種類 & 例 \\
大きさだけで決まる量 & 長さ・道のり、時間、質量、温度、体積、面積、速さ、値段 \\
向きと大きさで決まる量 & 移動（どこからどこへ）、速度、力 \\
\end{bktable}
\bktabcaption{「速さ」と「速度」は、別の量として区別する。}
\end{bkcenter}

\textbf{速さと速度}\quad 「速さ\qty{20}{m/s}」は、どっちへ走っているかを言っていません。「北へ\qty{20}{m/s}」のように、向きまでふくめたものを、物理では\textbf{速度}とよびます。速さは、速度の「大きさ」だけを取り出したものです。北へ\qty{20}{m/s}の車と、南へ\qty{20}{m/s}の車は、速さは同じですが、速度はちがいます。

このように、向きと大きさをもつ量を\textbf{ベクトル}（vector）といいます。\bkidx{べくとる}{ベクトル}大きさだけで決まる量を\textbf{スカラー}（scalar）といいます。\bkidx{すからー}{スカラー}ベクトルはラテン語の「運ぶもの」ということばから、スカラーは「はしご・目もり」ということばから来ています。スカラーは、1本の目もりの上の数で言いつくせる量、という意味です。

\begin{seikaku}
\textbf{数直線の上なら、向きは符号で表せた}\par
1本の道（数直線）の上では、向きは「＋の向き」と「－の向き」の2つしかありません。だから、2-2では向きを符号で表せました。東へ\qty{3}{m}は $+3$、西へ\qty{3}{m}は $-3$ です。平面の上では、向きは東・北東・北・北北西…と、いくらでもあります。符号の2つだけでは足りないので、矢印を使うのです。
\end{seikaku}

\bksubsection{2}{ベクトルを矢印と記号で表す}

\label{bk:4-10:kigou}ベクトルは、矢印でかきます。\textbf{矢印の向きがベクトルの向き}、\textbf{矢印の長さがベクトルの大きさ}です。力なら、「\qty{10}{N}を\qty{1}{cm}」のように、長さの縮尺を決めてかきます。

\noindent\begin{minipage}[c]{0.48\linewidth}\centering
\begin{tikzpicture}[x=1mm, y=1mm]
  \draw[kivk vec, bkBlue] (0,0) -- (20:40);
  \fill[bkInk] (0,0) circle (1);
  \node[font=\sffamily\small, anchor=north east] at (-0.8,-0.5) {A};
  \node[font=\sffamily\small, anchor=south west] at ($(20:40)+(0.5,0.5)$) {B};
  \node[font=\sffamily\footnotesize, anchor=north, text=bkGray] at (0,-6) {始点};
  \node[font=\sffamily\footnotesize, anchor=south, text=bkGray] at ($(20:40)+(0,6)$) {終点};
  % 長さの寸法線（矢印から 7 mm はなす）
  \draw[bkGray, line width=0.5pt, {Bar[width=3pt]}-{Bar[width=3pt]}] ($(0,0)+(110:7)$) -- ($(20:40)+(110:7)$);
  \node[font=\sffamily\footnotesize, rotate=20, anchor=south, fill=white, inner sep=0.8pt] at ($(20:20)+(110:7.6)$) {長さ＝大きさ（\qty{40}{N}）};
  \draw[bkGray, line width=0.4pt, -{Stealth[length=3pt]}] ($(20:30)+(-70:6)$) -- ($(20:38)+(-70:6)$);
  \node[font=\sffamily\footnotesize, anchor=north, text=bkGray] at ($(20:31)+(-70:7)$) {向き};
\end{tikzpicture}
\end{minipage}\hfill
\begin{minipage}[c]{0.49\linewidth}\raggedright 
矢印のもと（A）を\textbf{始点}、先（B）を\textbf{終点}といいます。\par\vspace{2pt}
左の図は、「\qty{10}{N}を\qty{1}{cm}」の縮尺で、\qty{40}{N}の力を長さ\qty{4}{cm}の矢印にかいたものです。
\end{minipage}

\par\vspace{4pt}
ベクトルは、次の記号で書きます。

\begin{itemize}[leftmargin=1.2\zw, itemsep=3pt, topsep=2pt, label={\color{\bkTT}●}]
\item \textbf{始点と終点で書く}\quad Aから始まりBで終わる矢印を、$\overrightarrow{\mathrm{AB}}$ と書き、「ベクトル・エー・ビー」と読みます。文字の順番が大事です。$\overrightarrow{\mathrm{BA}}$ は、Bから始まりAで終わる、反対向きの矢印です。
\item \textbf{1つの文字で書く}\quad $\vec{a}$、$\vec{b}$、$\vec{F}$（力、force の F）、$\vec{v}$（速度、velocity の v）のように、文字の上に小さな矢印をのせて書きます。$\vec{a}$ は「ベクトル・エー」と読みます。
\item \textbf{大きさ}\quad ベクトル $\vec{a}$ の大きさ（矢印の長さ）を $|\vec{a}|$ と書きます。$|\overrightarrow{\mathrm{AB}}|$ は、線分ABの長さです。
\end{itemize}

$|\ |$ は、3-2の絶対値と同じ記号です。絶対値 $|-3|=3$ は「数直線の上で、0からどれだけはなれているか」でした。$|\vec{a}|$ も「矢印のもとから先まで、どれだけはなれているか」です。どちらも、向き（符号）をすてて、大きさだけを取り出します。だから、$|\vec{a}|$ はただの数（スカラー）です。たとえば、$\vec{F}$ が「東へ\qty{40}{N}」の力なら、$|\vec{F}|=\qty{40}{N}$ です。

\bkivkmidashi{方眼の上の矢印の大きさ}
方眼の上にかいた矢印なら、大きさは三平方の定理（4-4）で求められます。下の図の $\vec{p}$ は、始点から右へ4マス、上へ3マス進んだ所が終点です。右へ4と上へ3を2辺とする直角三角形の斜辺なので、
\bkivkline{$|\vec{p}|=\sqrt{4^2+3^2}=\sqrt{16+9}=\sqrt{25}=5$}
1マスが\qty{1}{cm}なら、$\vec{p}$ の長さは\qty{5}{cm}です。

\bkivkmidashi{向きと大きさが同じなら、同じベクトル}

\begin{bkfigure}
\begin{tikzpicture}[x=5.6mm, y=5.6mm]
  \draw[bkLine, line width=0.3pt] (0,0) grid (20,7);
  \draw[kivk vec, bkBlue] (1,1) -- (5,4);
  \draw[kivk vec, bkBlue] (6,2) -- (10,5);
  \draw[kivk vec, bkBlue] (11,0) -- (15,3);
  \draw[kivk vec, bkOrange] (16,1) -- (19,5);
  \draw[kivk vec, bkGreen!70!black] (12,4.5) -- (14,6);
  \node[kivk lab, fill=white, text=bkBlue!80!black] at (2.4,3.3) {$\vec{p}$};
  \node[kivk lab, fill=white, text=bkBlue!80!black] at (7.4,4.3) {$\vec{q}$};
  \node[kivk lab, fill=white, text=bkBlue!80!black] at (12.4,2.3) {$\vec{r}$};
  \node[kivk lab, fill=white, text=bkOrange!80!black] at (16.8,3.6) {$\vec{s}$};
  \node[kivk lab, fill=white, text=bkGreen!50!black] at (12.4,5.8) {$\vec{t}$};
\end{tikzpicture}
\bkfigcaption{$\vec{p}$・$\vec{q}$・$\vec{r}$ は、どれも「右へ4・上へ3」で同じベクトル。$\vec{s}$ は大きさ5で同じだが向きがちがう。$\vec{t}$ は向きが同じだが大きさが半分。}\label{bk:4-10:fig-onaji}
\end{bkfigure}

図\ref{bk:4-10:fig-onaji}の $\vec{p}$、$\vec{q}$、$\vec{r}$ は、かいてある場所はちがいますが、どれも「右へ4・上へ3」の矢印です。3-2で、ゲームのキャラクターが「右へ4・上へ3」動くとき、どこから動いても、その「動き」は同じでした。ベクトルも同じで、\textbf{向きと大きさが同じなら、どこにかいても同じベクトル}と考えます。$\vec{p}=\vec{q}=\vec{r}$ と書きます。

$\vec{s}$ は、大きさは $\sqrt{3^2+4^2}=5$ で $\vec{p}$ と同じですが、向きがちがうので、$\vec{s}\neq\vec{p}$ です。$\vec{t}$ は、向きは同じですが、大きさが2.5で半分なので、$\vec{t}\neq\vec{p}$ です。2つのベクトルが等しいのは、向きと大きさの\textbf{両方}が同じときだけです。

\begin{seikaku}
\textbf{力の矢印は、どこにかいてもよい？}\par
ベクトルを計算するときは、矢印を平行に動かしてかき直してかまいません。ただし、力が物のどこにはたらくか（\textbf{作用点}）は、物を回す働きを考えるときに大事になります。ドアのノブを押すのと、ちょうつがいの近くを押すのとでは、同じ力でもドアの回り方がちがいます。これは4-14で調べます。この章では、物を回すことは考えず、物全体がどっちへ動くかだけを考えます。
\end{seikaku}

\bkkigocards{\bkkigocard{$\vec{a}$}{ベクトル・エー}{向きと大きさをもつ量。$\overrightarrow{\mathrm{AB}}$ はAからBへの矢印。}{vector}{4-10}}
{\bkkigocard{$|\vec{a}|$}{ベクトル・エーの大きさ}{矢印の長さ。向きをすてた数（スカラー）。}{magnitude}{4-10}}

\bksubsection{3}{ベクトルのたし算：矢印をつなぐ}

\label{bk:4-10:wa}\bkivkstep{1}{移動を続けると}
ソラさんは、家から自転車で東へ\qty{4}{km}走り、そこから北へ\qty{3}{km}走りました。いま、ソラさんは家から見て、どこにいるでしょう。

\noindent\begin{minipage}[c]{0.5\linewidth}\centering
\begin{tikzpicture}[x=9mm, y=9mm]
  \draw[bkLine, line width=0.3pt] (-0.5,-0.5) grid (5,3.6);
  % 家
  \filldraw[fill=bkYellow!50, draw=bkYellow!50!black, line width=0.5pt] (-0.3,-0.25) rectangle (0.3,0.2);
  \filldraw[fill=bkRed!60, draw=bkRed!70!black, line width=0.5pt] (-0.38,0.2) -- (0,0.5) -- (0.38,0.2) -- cycle;
  \draw[kivk vec, bkBlue] (0.4,0) -- (4,0);
  \draw[kivk vec, bkOrange] (4,0.08) -- (4,3);
  \draw[kivk vec, bkRed] (0.32,0.24) -- (3.92,2.94);
  \fill[bkInk] (4,3) circle (0.09);
  \node[kivk lab, text=bkBlue!80!black, anchor=north, fill=white] at (2.2,-0.12) {$\vec{a}$　東へ4 km};
  \node[kivk lab, text=bkOrange!80!black, anchor=west, fill=white] at (4.15,1.5) {$\vec{b}$　北へ3 km};
  \node[kivk lab, text=bkRed!80!black, anchor=south east, fill=white] at (1.7,1.75) {$\vec{a}+\vec{b}$};
\end{tikzpicture}
\end{minipage}\hfill
\begin{minipage}[c]{0.47\linewidth}\raggedright 
1マスは\qty{1}{km}です。\par\vspace{2pt}
東へ\qty{4}{km}の移動を $\vec{a}$（青）、北へ\qty{3}{km}の移動を $\vec{b}$（オレンジ）とします。\par\vspace{2pt}
2つの移動を続けた結果は、家から終わりの場所への1本の矢印（赤）です。
\end{minipage}

\par\vspace{4pt}
赤い矢印は、直角をはさむ2辺が4と3の直角三角形の斜辺なので、長さは $\sqrt{4^2+3^2}=5$ です。ソラさんは、家から北東よりの向きへ、まっすぐはかって\qty{5}{km}の所にいます。走った道のりは $4+3=\qty{7}{km}$ ですが、家からの移動は\qty{5}{km}です。

このように、1つめのベクトル $\vec{a}$ の終点に、2つめのベクトル $\vec{b}$ の始点をつなぎ、$\vec{a}$ の始点から $\vec{b}$ の終点へ引いた矢印を、$\vec{a}$ と $\vec{b}$ の\textbf{和}といい、$\vec{a}+\vec{b}$ と書きます。\bkidx{べくとるのわ}{ベクトルの和}「$\vec{a}$ だけ動いて、続けて $\vec{b}$ だけ動くのは、$\vec{a}+\vec{b}$ だけ動くのと同じ」という意味です。

始点と終点の記号で書くと、とてもわかりやすくなります。
\bkivkline{\fbox{$\overrightarrow{\mathrm{AB}}+\overrightarrow{\mathrm{BC}}=\overrightarrow{\mathrm{AC}}$}}
「AからBへ行き、続けてBからCへ行くのは、AからCへ行くのと同じ」です。まん中のBが、つなぎめになって消えます。

\bkivkstep{2}{2つの力を合わせると：平行四辺形}
では、\bkhatenaref のように、1つの荷物を2つの力が同時に引くときはどうでしょう。2つの力は、同じ点から出る2本の矢印です。つなぐ形にはなっていません。

やってみようでは、2つの力の矢印を2辺とする平行四辺形をかくと、その\textbf{対角線}が、2つの力を合わせた働きになりました。この「力の平行四辺形」は、数学で決めたことではなく、実験で何度も確かめられてきた自然のきまり（法則、4-1）です。

\begin{bkfigure}
\begin{tikzpicture}[x=1mm, y=1mm]
  % 左：平行四辺形
  \draw[bkGray, line width=0.5pt, dashed] (32,0) -- (32,24) -- (0,24);
  \draw[kivk vec, bkBlue] (0,0) -- (32,0);
  \draw[kivk vec, bkOrange] (0,0) -- (0,24);
  \draw[kivk vec, bkRed] (0.8,0.6) -- (31.2,23.4);
  \fill[bkInk] (0,0) circle (1);
  \node[kivk lab, text=bkBlue!80!black, anchor=north] at (16,-1.5) {$\vec{a}$};
  \node[kivk lab, text=bkOrange!80!black, anchor=east] at (-1.5,12) {$\vec{b}$};
  \node[kivk lab, text=bkRed!80!black, anchor=south east, fill=white] at (12,12.5) {$\vec{a}+\vec{b}$};
  \node[font=\sffamily\small, anchor=north] at (16,-8) {平行四辺形の対角線};
  % 右：つなぐ
  \begin{scope}[shift={(62,0)}]
  \draw[bkGray, line width=0.5pt, dashed] (0,0.6) -- (0,24) -- (31.4,24);
  \draw[kivk vec, bkBlue] (0,0) -- (32,0);
  \draw[kivk vec, bkOrange] (32,0.2) -- (32,24);
  \draw[kivk vec, bkRed] (0.8,0.6) -- (31.2,23.4);
  \fill[bkInk] (0,0) circle (1);
  \node[kivk lab, text=bkBlue!80!black, anchor=north] at (16,-1.5) {$\vec{a}$};
  \node[kivk lab, text=bkOrange!80!black, anchor=west] at (33.5,12) {$\vec{b}$};
  \node[kivk lab, text=bkRed!80!black, anchor=south east, fill=white] at (12,12.5) {$\vec{a}+\vec{b}$};
  \node[font=\sffamily\small, anchor=north] at (16,-8) {$\vec{b}$ を平行に動かしてつなぐ};
  \end{scope}
\end{tikzpicture}
\bkfigcaption{平行四辺形の向かい合う辺は、平行で長さが等しい。だから右がわの辺は、$\vec{b}$ と同じベクトル。}\label{bk:4-10:fig-heikoushihen}
\end{bkfigure}

図\ref{bk:4-10:fig-heikoushihen}の左の平行四辺形の右がわの辺を見てください。平行四辺形の向かい合う辺は、平行で長さが等しいので、この辺は $\vec{b}$ と同じ向き・同じ大きさです。つまり、$\vec{b}$ を平行に動かして $\vec{a}$ の先につないだ矢印です（右の図）。だから、\textbf{平行四辺形の対角線は、矢印をつないだ和と同じ}です。2つの方法は、同じ答えを出します。

2つ以上の力を合わせて1つにした力を、\textbf{合力}（ごうりょく）といいます。\bkidx{ごうりょく}{合力}力の合力は、ベクトルの和です。

\bkivkmidashi{\bkhatenaref の答え}
ソラさんの力を $\vec{a}$（東へ\qty{40}{N}）、ミナさんの力を $\vec{b}$（北へ\qty{30}{N}）とします。合力 $\vec{a}+\vec{b}$ は、2辺が40と30の長方形の対角線なので、大きさは
\bkivkline{$|\vec{a}+\vec{b}|=\sqrt{40^2+30^2}=\sqrt{1600+900}=\sqrt{2500}=50$}
で\qty{50}{N}です。向きは、東から北へ角 $\theta$ だけ回った向きで、$\tan\theta=\dfrac{30}{40}=0.75$ です（4-6）。電卓で $\tan^{-1}0.75$ を計算すると、$\theta=36.86\cdots°$ です（電卓は DEG にしておく）。

荷物は、\textbf{東から約37°北へ回った向き}に、\qty{50}{N}の力1つで引かれたのと同じように動き出します。床との摩擦は、動きをじゃまする向きにはたらくので、動き出すかどうかは変えますが、動き出す向きは変えません。

\bkng{2人の力を合わせると $40+30=\qty{70}{N}$}
\noindent 大きさをそのままたしてよいのは、2つの力が\textbf{同じ向き}のときだけです。

\bkivkstep{3}{同じ向き・反対向きの和}
一直線上の2つの力の和を、3つの場合でくらべます。

\begin{bkfigure}
\begin{tikzpicture}[x=1mm, y=1mm]
  \foreach \y/\t in {58/同じ向き, 44/反対向き, 12/直角} {\node[font=\sffamily\small, anchor=east] at (-3,\y) {\t};}
  % 同じ向き
  \draw[kivk vec, bkBlue] (0,58) -- (40,58);
  \draw[kivk vec, bkOrange] (40.5,58) -- (70,58);
  \node[kivk lab, text=bkBlue!80!black, anchor=south] at (20,59.5) {\qty{40}{N}};
  \node[kivk lab, text=bkOrange!80!black, anchor=south] at (55,59.5) {\qty{30}{N}};
  \node[kivk lab, text=bkRed!80!black, anchor=west] at (73,58) {合力\qty{70}{N}};
  % 反対向き
  \draw[kivk vec, bkBlue] (0,45) -- (40,45);
  \draw[kivk vec, bkOrange] (40,43) -- (10.5,43);
  \draw[kivk vec, bkRed, line width=1.2pt] (0,40.2) -- (10,40.2);
  \node[kivk lab, text=bkBlue!80!black, anchor=south] at (20,46.5) {\qty{40}{N}};
  \node[kivk lab, text=bkOrange!80!black, anchor=north] at (28,41.6) {\qty{30}{N}};
  \node[kivk lab, text=bkRed!80!black, anchor=west] at (73,44) {合力\qty{10}{N}（$40-30$）};
  % 直角
  \draw[kivk vec, bkBlue] (0,0) -- (32,0);
  \draw[kivk vec, bkOrange] (32,0.2) -- (32,24);
  \draw[kivk vec, bkRed] (0.6,0.6) -- (31.4,23.4);
  \node[kivk lab, text=bkBlue!80!black, anchor=north] at (16,-1.5) {\qty{40}{N}};
  \node[kivk lab, text=bkOrange!80!black, anchor=west] at (33.5,12) {\qty{30}{N}};
  \node[kivk lab, text=bkRed!80!black, anchor=west] at (73,12) {合力\qty{50}{N}（$\sqrt{40^2+30^2}$）};
\end{tikzpicture}
\bkfigcaption{同じ\qty{40}{N}と\qty{30}{N}でも、向きによって合力は\qty{10}{N}から\qty{70}{N}まで変わる（直角の図は縮尺を0.8倍にした）。}
\end{bkfigure}

反対向きのときは、東を＋とすると $(+40)+(-30)=+10$ で、東へ\qty{10}{N}です（2-2）。綱引きで、強いほうへ少しずつ動いていくのと同じです。

三角形の2辺の長さの和は、残りの1辺より長くなります。だから、合力の大きさは、いつも「2つの力の大きさの和」以下です。和とちょうど同じになるのは、同じ向きのときだけです。
\bkivkline{$|\vec{a}+\vec{b}|\leqq|\vec{a}|+|\vec{b}|$}

\bkivkstep{4}{たす順番を入れかえても同じ}
図\ref{bk:4-10:fig-heikoushihen}の平行四辺形で、左下から右上の角へ行く道は2通りあります。下の辺から右の辺へ行く道は $\vec{a}+\vec{b}$、左の辺から上の辺へ行く道は $\vec{b}+\vec{a}$ です。どちらも同じ対角線に着くので、
\bkivkline{$\vec{a}+\vec{b}=\vec{b}+\vec{a}$}
です。数のたし算と同じように、たす順番を入れかえてもかまいません。3つ以上のベクトルの和は、矢印を次々につないでいけば求められます。

\bkivkstep{5}{角があるときの合力の大きさ：余弦定理}
2つの力が直角でないときは、平行四辺形の対角線の長さを、余弦定理（4-7）で求めます。\qty{5}{N}の力 $\vec{a}$ と\qty{3}{N}の力 $\vec{b}$ が、60°の角をなしているとします。

\noindent\begin{minipage}[c]{0.5\linewidth}\centering
\begin{tikzpicture}[x=1mm, y=1mm]
  \draw[bkGray, line width=0.5pt, dashed] (12,20.785) -- (52,20.785);
  \draw[kivk vec, bkBlue] (0,0) -- (40,0);
  \draw[kivk vec, bkOrange] (0,0) -- (12,20.785);
  \draw[kivk vec, bkOrange!60] (40,0) -- (52,20.785);
  \draw[kivk vec, bkRed] (0.9,0.36) -- (51.1,20.4);
  \fill[bkInk] (0,0) circle (1);
  \draw[bkInk, line width=0.5pt] (7,0) arc (0:60:7);
  \node[font=\sffamily\footnotesize] at (30:10.5) {60°};
  \draw[bkInk, line width=0.5pt] ($(40,0)+(60:6)$) arc (60:180:6);
  \node[font=\sffamily\footnotesize, fill=white, inner sep=0.5pt] at ($(40,0)+(120:9.6)$) {120°};
  \node[kivk lab, text=bkBlue!80!black, anchor=north] at (20,-1.5) {$|\vec{a}|=5$};
  \node[kivk lab, text=bkOrange!80!black, anchor=east] at (5,12) {$|\vec{b}|=3$};
  \node[kivk lab, text=bkOrange!70!black, anchor=west] at (47.5,9) {3};
  \node[kivk lab, text=bkRed!80!black, anchor=south, fill=white] at (26,13) {$\vec{a}+\vec{b}$};
\end{tikzpicture}
\end{minipage}\hfill
\begin{minipage}[c]{0.47\linewidth}\raggedright 
$\vec{b}$ を平行に動かして $\vec{a}$ の先につなぐと（うすいオレンジ）、赤い対角線をふくむ三角形ができます。\par\vspace{2pt}
この三角形の、$\vec{a}$ の先の角は120°です。平行四辺形のとなり合う角の和は180°だからです（平行線と角、4-1）。
\end{minipage}

\par\vspace{4pt}
三角形の2辺5と3と、そのあいだの角120°がわかったので、余弦定理で向かいの辺（合力の大きさ）を求めます。$\cos120°=-0.5$ です。
\begin{align*}
|\vec{a}+\vec{b}|^2&=5^2+3^2-2\times5\times3\times\cos120°\\
&=25+9-30\times(-0.5)=25+9+15=49
\end{align*}
$|\vec{a}+\vec{b}|=\sqrt{49}=7$ で、合力は\qty{7}{N}です。$5+3=8$ より小さく、直角のときの $\sqrt{5^2+3^2}=\sqrt{34}$ より大きくなりました。

2つの力のなす角を $\theta$ としたときの合力の大きさを、表にします。（$\approx$ は「ほぼ等しい」という記号です。5-8でくわしく学びます。）

\begin{bkcenter}
\begin{bktable}{colspec={Q[l,m,wd=24mm] Q[c,m,wd=10mm] Q[c,m,wd=10mm] Q[c,m,wd=25mm] Q[c,m,wd=25mm] Q[c,m,wd=10mm]}, rowsep=1.5pt, colsep=4pt}
なす角 $\theta$ & 0° & 60° & 90° & 120° & 180° \\
合力の大きさ（N） & 8 & 7 & $\sqrt{34}\approx5.83$ & $\sqrt{19}\approx4.36$ & 2 \\
\end{bktable}
\bktabcaption{\qty{5}{N}と\qty{3}{N}の合力。2つの力の向きが開くほど、合力は小さくなる。}
\end{bkcenter}

\Needspace{16\baselineskip}
\begin{motto}[3]
\textbf{3つの力のつり合いは、三角形で解ける}\par
\noindent\begin{minipage}[c]{0.6\linewidth}
重さ\qty{10}{N}の荷物を、2本のロープでつるします。2本のロープは直角に開いていて、たての線から30°と60°かたむいています。荷物が止まっているとき、2本のロープの力と重さの3つの力はつり合っているので、矢印をつなぐと閉じた三角形になります（④で学びます）。この三角形は、斜辺が10（重さ）、角が30°と60°の直角三角形です。だから、ロープの力は $10\cos30°\approx\qty{8.7}{N}$ と $10\cos60°=\qty{5}{N}$ です（4-6）。たてに近いロープのほうが、大きな力で支えています。直角三角形にならないときは、4-7の正弦定理を使って、同じように求められます。
\end{minipage}\hfill
\begin{minipage}[c]{0.36\linewidth}\centering
\begin{tikzpicture}[x=1mm, y=1mm]
  \draw[kivk vec, bkPurple, line width=1.2pt] (0,30) -- (0,0.6);
  \draw[kivk vec, bkBlue, line width=1.2pt] (0,0) -- (12.99,7.5);
  \draw[kivk vec, bkOrange, line width=1.2pt] (12.99,7.5) -- (0.4,29.4);
  \draw[bkInk, line width=0.4pt] ($(12.99,7.5)+(210:2.5)$) -- ++(120:2.5) -- ++(30:2.5);
  \node[kivk lab, text=bkPurple!80!black, anchor=east] at (-1.2,15) {重さ10};
  \node[kivk lab, text=bkBlue!80!black, anchor=north west] at (7,3) {5};
  \node[kivk lab, text=bkOrange!80!black, anchor=south west] at (7.5,19.5) {8.7};
\end{tikzpicture}
\end{minipage}
\end{motto}

\bksubsection{4}{実数倍・反対のベクトル・ゼロベクトル}

\label{bk:4-10:jissubai}東へ\qty{5}{m/s}で走る自転車が、同じ向きに2倍の速さで走ると、速度は東へ\qty{10}{m/s}です。矢印でいうと、同じ向きで長さが2倍です。これは、同じ矢印を2本つないだ $\vec{a}+\vec{a}$ でもあるので、$2\vec{a}$ と書きます。

\begin{bkfigure}
\begin{tikzpicture}[x=1mm, y=1mm]
  \draw[kivk vec, bkBlue] (0,0) -- ++(25:16);
  \draw[kivk vec, bkBlue] (22,0) -- ++(25:32);
  \draw[kivk vec, bkBlue] (56,0) -- ++(25:8);
  \draw[kivk vec, bkRed] ($(84,0)+(25:16)$) -- (84,0);
  \draw[kivk vec, bkRed] ($(102,0)+(25:32)$) -- (102,0);
  \node[kivk lab, anchor=north] at (7,-3) {$\vec{a}$};
  \node[kivk lab, anchor=north] at (36,-3) {$2\vec{a}$};
  \node[kivk lab, anchor=north] at (60,-3) {$0.5\vec{a}$};
  \node[kivk lab, anchor=north] at (91,-3) {$-\vec{a}$};
  \node[kivk lab, anchor=north] at (116,-3) {$-2\vec{a}$};
\end{tikzpicture}
\bkfigcaption{正の数をかけると同じ向き、負の数をかけると反対向き。長さは、かけた数の絶対値の倍。}
\end{bkfigure}

ふつうの数（0、正の数、負の数、分数、小数、$\sqrt2$ など。まとめて\textbf{実数}（じっすう）といいます）を $k$ とします。ベクトル $\vec{a}$ を $k$ 倍した $k\vec{a}$ を、次のように決めます（これは定義です。4-1）。

\begin{itemize}[leftmargin=1.2\zw, itemsep=2pt, topsep=2pt, label={\color{\bkTT}●}]
\item $k$ が正のとき：$\vec{a}$ と同じ向きで、長さが $k$ 倍の矢印。
\item $k$ が負のとき：$\vec{a}$ と\textbf{反対の向き}で、長さが $|k|$ 倍の矢印。
\item $k=0$ のとき：長さ0の「矢印」。これを\textbf{ゼロベクトル}といい、$\vec{0}$ と書きます。
\end{itemize}

負の数をかけると向きが反対になるのは、2-2の「マイナスは反対の向き」と同じ考えです。数直線の上で $\times(-1)$ をすると、0をはさんで反対がわへ移りました。とくに、$(-1)\vec{a}$ を $-\vec{a}$ と書き、$\vec{a}$ の\textbf{反対のベクトル}といいます。$\vec{a}$ と大きさが同じで、向きが反対です。$\overrightarrow{\mathrm{AB}}$ の反対のベクトルは $\overrightarrow{\mathrm{BA}}$ です。

\bkivkmidashi{ゼロベクトル}
東へ\qty{3}{m}歩いてから、西へ\qty{3}{m}歩くと、もとの場所にもどります。結果として、どこへも動いていません。
\bkivkline{$\vec{a}+(-\vec{a})=\vec{0}$}
ゼロベクトル $\vec{0}$ は、大きさが0のベクトルです。長さがないので、向きは決まりません。数の0と同じように、どんなベクトルにたしても、相手を変えません（$\vec{a}+\vec{0}=\vec{a}$）。

やってみようの3つの力A、B、Cを、矢印を次々につないでみましょう。AとBの和はCと同じ長さで反対向きなので、Cをつなぐと、出発した点にもどってきます。3つの力の和は $\vec{0}$ です。\textbf{力がつり合っている}とは、\textbf{合力が $\vec{0}$} になっていることです。つないだ矢印は、閉じた三角形になります。

\begin{bkfigure}
\begin{tikzpicture}[x=1mm, y=1mm]
  \draw[kivk vec, bkBlue] (0,0) -- (24,0);
  \draw[kivk vec, bkOrange] (24,0.2) -- (24,32);
  \draw[kivk vec, bkGreen!70!black] (23.4,31.4) -- (0.5,0.7);
  \node[kivk lab, text=bkBlue!80!black, anchor=north] at (12,-1.5) {A（\qty{3}{N}）};
  \node[kivk lab, text=bkOrange!80!black, anchor=west] at (25.5,16) {B（\qty{4}{N}）};
  \node[kivk lab, text=bkGreen!50!black, anchor=south east] at (10.5,17) {C（\qty{5}{N}）};
  \node[font=\sffamily\small, anchor=west, align=left, text width=62mm] at (44,16) {A、B、Cをつなぐと、\\三角形が閉じて出発点にもどる。\\$\text{A}+\text{B}+\text{C}=\vec{0}$（つり合い）};
\end{tikzpicture}
\bkfigcaption{つり合っている3つの力の矢印をつなぐと、閉じた三角形になる（\qty{1}{N}を\qty{8}{mm}でかいた）。}
\end{bkfigure}

\bkivkmidashi{平行なベクトル}
$\vec{0}$ でない2つのベクトル $\vec{a}$、$\vec{b}$ が、同じ向きか反対の向きのとき、$\vec{a}$ と $\vec{b}$ は\textbf{平行}であるといいます。このとき、$\vec{b}=k\vec{a}$ となる実数 $k$ があります。たとえば、図の $2\vec{a}$ も $-\vec{a}$ も、$\vec{a}$ と平行です。

\bksubsection{5}{ベクトルのひき算：差と、速度の変化 $\Delta\vec{v}$}

\label{bk:4-10:sa}\bkivkstep{1}{ひくことは、反対をたすこと}
2-2では、「ひくことは、反対の数をたすこと」でした。$5-3=5+(-3)$ です。ベクトルでも同じように決めます。
\bkivkline{\fbox{$\vec{a}-\vec{b}=\vec{a}+(-\vec{b})$}}
$\vec{b}$ をひくことは、反対のベクトル $-\vec{b}$ をたすことです。

\begin{bkfigure}
\begin{tikzpicture}[x=1mm, y=1mm]
  % 左：a+(-b)
  \draw[kivk vec, bkBlue] (0,0) -- (36,8);
  \draw[kivk vec, bkOrange!60, dashed] (0,0) -- (12,24);
  \draw[kivk vec, bkOrange] (36,7.6) -- (24.2,-15.6);
  \draw[kivk vec, bkRed] (0.8,-0.5) -- (23.4,-15.4);
  \fill[bkInk] (0,0) circle (1);
  \node[font=\sffamily\footnotesize, anchor=east] at (-1.5,0) {O};
  \node[kivk lab, text=bkBlue!80!black, anchor=south] at (22,6.8) {$\vec{a}$};
  \node[kivk lab, text=bkOrange!70!black, anchor=east] at (4.5,13) {$\vec{b}$};
  \node[kivk lab, text=bkOrange!80!black, anchor=west] at (31.5,-5) {$-\vec{b}$};
  \node[kivk lab, text=bkRed!80!black, anchor=north east] at (10,-10) {$\vec{a}+(-\vec{b})$};
  \node[font=\sffamily\small, anchor=north] at (18,-21) {$\vec{a}$ に $-\vec{b}$ をつなぐ};
  % 右：b の先から a の先へ
  \begin{scope}[shift={(66,-6)}]
  \draw[kivk vec, bkBlue] (0,0) -- (36,8);
  \draw[kivk vec, bkOrange] (0,0) -- (12,24);
  \draw[kivk vec, bkRed] (12.5,23.1) -- (35.5,8.8);
  \fill[bkInk] (0,0) circle (1);
  \node[font=\sffamily\footnotesize, anchor=east] at (-1.5,0) {O};
  \node[font=\sffamily\footnotesize, anchor=north west] at (36.5,7.5) {A};
  \node[font=\sffamily\footnotesize, anchor=south] at (12,25) {B};
  \node[kivk lab, text=bkBlue!80!black, anchor=north] at (20,2.6) {$\vec{a}$};
  \node[kivk lab, text=bkOrange!80!black, anchor=east] at (4.5,13) {$\vec{b}$};
  \node[kivk lab, text=bkRed!80!black, anchor=south west] at (26,17.5) {$\vec{a}-\vec{b}$};
  \node[font=\sffamily\small, anchor=north] at (18,-15) {$\vec{b}$ の先から $\vec{a}$ の先へ};
  \end{scope}
\end{tikzpicture}
\bkfigcaption{2つの赤い矢印は、向きも大きさも同じ。$\vec{a}-\vec{b}$ は、始点をそろえたとき「$\vec{b}$ の先から $\vec{a}$ の先へ」の矢印。}\label{bk:4-10:fig-sa}
\end{bkfigure}

図\ref{bk:4-10:fig-sa}の左は、$\vec{a}$ の先に $-\vec{b}$ をつないで、$\vec{a}+(-\vec{b})$ を作ったものです。右は、$\vec{a}$ と $\vec{b}$ の始点をOにそろえて、$\vec{b}$ の先Bから $\vec{a}$ の先Aへ矢印を引いたものです。2つの赤い矢印は、向きも大きさも同じです。

\bkivkstep{2}{もう1つの見方：「何をたせば、そうなるか」}
数のひき算 $5-3$ は、「3に何をたせば5になるか」の答えでした（$3+2=5$ なので2）。ベクトルでも、$\vec{a}-\vec{b}$ は、\textbf{$\vec{b}$ に何をたせば $\vec{a}$ になるか}の答えです。図\ref{bk:4-10:fig-sa}の右で、Oから $\vec{b}$ だけ進んでBに着き、そこから赤い矢印だけ進むと、Aに着きます。$\vec{b}+(\text{赤い矢印})=\vec{a}$ なので、赤い矢印が $\vec{a}-\vec{b}$ です。

始点と終点の記号で書くと、
\bkivkline{\fbox{$\overrightarrow{\mathrm{OA}}-\overrightarrow{\mathrm{OB}}=\overrightarrow{\mathrm{BA}}$}}
「あと（A）－まえ（B）」で、BからAへの矢印です。3-2で、どれだけ動いたかを「あと－まえ」で求めたのと同じです。ひく順番を逆にすると、$\overrightarrow{\mathrm{OB}}-\overrightarrow{\mathrm{OA}}=\overrightarrow{\mathrm{AB}}$ で、向きが反対の矢印になります。

\bkng{$\vec{a}-\vec{b}$ は、$\vec{a}$ の先から $\vec{b}$ の先へ}
\bkok{$\vec{a}-\vec{b}$ は、$\vec{b}$ の先から $\vec{a}$ の先へ（ひくほうから、ひかれるほうへ）}

\bkivkstep{3}{速度の変化 $\Delta\vec{v}$：一直線上}
2-3で、変化量は「あと－まえ」でした。速度の変化も同じで、\textbf{速度の変化＝あとの速度－まえの速度}です。速度はベクトルなので、変化もベクトルの差で求めます。
\bkivkline{$\Delta\vec{v}=\vec{v}_{\text{あと}}-\vec{v}_{\text{まえ}}$}
$\Delta$ は「変化量」の記号です（2-3）。

\noindent\begin{minipage}[c]{0.5\linewidth}\centering
\begin{tikzpicture}[x=1mm, y=1mm]
  \fill[bkGray!35] (46,-16) rectangle (50,14);
  \draw[bkGray, line width=0.7pt] (46,-16) -- (46,14);
  \foreach \y in {-14,-10,...,12} {\draw[bkGray, line width=0.4pt] (46,\y) -- (49,\y+2.5);}
  \fill[bkBlue!60] (41,9) circle (2.6);
  \draw[kivk vec, bkBlue] (21,9) -- (37,9);
  \node[kivk lab, text=bkBlue!80!black, anchor=south] at (29,10.5) {まえ：東へ8};
  \fill[bkOrange!70] (41,-1) circle (2.6);
  \draw[kivk vec, bkOrange] (37,-1) -- (25,-1);
  \node[kivk lab, text=bkOrange!80!black, anchor=south] at (31,0.5) {あと：西へ6};
  \draw[kivk vec, bkRed] (37,-11) -- (9,-11);
  \node[kivk lab, text=bkRed!80!black, anchor=north] at (23,-12.5) {$\Delta\vec{v}$：西へ14};
  \draw[bkGray, line width=0.6pt, -{Stealth[length=3pt]}] (2,17) -- (14,17);
  \node[font=\sffamily\footnotesize, anchor=west] at (14.5,17) {東が＋};
\end{tikzpicture}
\end{minipage}\hfill
\begin{minipage}[c]{0.47\linewidth}\raggedright 
ボールが東へ\qty{8}{m/s}で壁に当たり、西へ\qty{6}{m/s}ではね返りました。東を＋とすると、まえの速度は $+8$、あとの速度は $-6$ です。
\[ \Delta v=(-6)-(+8)=-14 \]
速度の変化は、西へ\qty{14}{m/s}です。
\end{minipage}

\par\vspace{4pt}
速さだけを見ると、8から6へ「2へった」ように見えます。でも速度は、東へ8の状態から、まず0になり、さらに西へ6まで変わりました。変化は $8+6=14$ です。壁は、ボールの速度をこれだけ変える大きな力で、ボールを押し返したのです。

\bkivkstep{4}{速度の変化 $\Delta\vec{v}$：向きが変わるとき}
\label{bk:4-10:dv}車が、速さ\qty{10}{m/s}のまま、東向きから北向きへ、交差点を左に曲がりました。速さは変わっていません。では、速度は変わったでしょうか。

\begin{bkfigure}
\begin{tikzpicture}[x=1mm, y=1mm]
  % 左：道路と車の速度
  \fill[bkGray!18] (-6,-6) rectangle (34,6);
  \fill[bkGray!18] (22,-6) rectangle (34,32);
  \draw[bkGray, line width=0.6pt, dashed] (-6,0) -- (16,0) arc (-90:0:12) -- (28,32);
  \fill[bkInk] (4,0) circle (1);
  \draw[kivk vec, bkBlue] (4,0) -- (18,0);
  \node[kivk lab, text=bkBlue!80!black, anchor=north] at (8,-1.5) {$\vec{v}_{\text{まえ}}$};
  \fill[bkInk] (28,18) circle (1);
  \draw[kivk vec, bkOrange] (28,18) -- (28,32);
  \node[kivk lab, text=bkOrange!80!black, anchor=east] at (26.5,24) {$\vec{v}_{\text{あと}}$};
  \node[font=\sffamily\small, anchor=north] at (14,-10) {上から見た交差点};
  \bkivkcompass{(-6,22)}
  % 右：ベクトルの図
  \begin{scope}[shift={(66,-2)}]
  \draw[kivk vec, bkBlue] (0,0) -- (28,0);
  \draw[kivk vec, bkOrange] (0,0) -- (0,28);
  \draw[kivk vec, bkRed] (27.4,0.6) -- (0.6,27.4);
  \fill[bkInk] (0,0) circle (1);
  \node[kivk lab, text=bkBlue!80!black, anchor=north] at (14,-1.5) {東へ10};
  \node[kivk lab, text=bkOrange!80!black, anchor=east] at (-1.5,14) {北へ10};
  \node[kivk lab, text=bkRed!80!black, anchor=south west] at (14.5,14.5) {$\Delta\vec{v}$};
  \node[font=\sffamily\small, anchor=north] at (14,-8) {始点をそろえる};
  \end{scope}
\end{tikzpicture}
\bkfigcaption{始点をそろえると、$\Delta\vec{v}=\vec{v}_{\text{あと}}-\vec{v}_{\text{まえ}}$ は「まえの先からあとの先へ」の矢印（赤）。北西を向く。}\label{bk:4-10:fig-dv}
\end{bkfigure}

まえの速度は東へ\qty{10}{m/s}、あとの速度は北へ\qty{10}{m/s}です。始点をそろえてかくと、$\Delta\vec{v}$ は「まえの先から、あとの先へ」の矢印です（②）。この矢印は、直角をはさむ2辺が10と10の直角二等辺三角形の斜辺なので、
\bkivkline{$|\Delta\vec{v}|=\sqrt{10^2+10^2}=\sqrt{200}=10\sqrt2\approx14.1$}
です（$1:1:\sqrt2$、4-4）。速度は、北西の向きに\qty{14.1}{m/s}も変わりました。

もう1つの考え方でも確かめましょう。$\Delta\vec{v}=\vec{v}_{\text{あと}}+(-\vec{v}_{\text{まえ}})$ です。$-\vec{v}_{\text{まえ}}$ は西へ\qty{10}{m/s}なので、「北へ10」と「西へ10」の和になり、やはり北西へ $10\sqrt2$ です。

\textbf{速さが同じでも、向きが変われば速度は変わります}（$\Delta\vec{v}\neq\vec{0}$）。物理シリーズの力学では、速度を変えるには力がいることを学びます。カーブを曲がる車には、タイヤと道路のあいだの摩擦の力が、曲がる内がわへ向かってはたらいています。雪の日に道がすべりやすいと、この力がたりなくなって、車は曲がりきれずに外へ出てしまいます。

\begin{motto}[3]
\textbf{円を回る物の $\Delta\vec{v}$ は、円の中心を向く}\par
\noindent\begin{minipage}[c]{0.56\linewidth}\raggedright 
同じ速さで円を回る物の速度は、いつも円の接線の向きです（4-2）。円の上の点Pと、少し先の点Qで、速度の矢印をかきます（右の図の青とオレンジ）。始点をそろえて差 $\Delta\vec{v}$ をかくと（赤）、その向きは、PとQのまん中から\textbf{円の中心へ向かう}向きになります。\par
PとQを近づけても、$\Delta\vec{v}$ はいつも中心を向いています。物理では、この中心を向く速度の変化を生む力を、\textbf{向心力}（こうしんりょく）とよびます。
\end{minipage}\hfill
\begin{minipage}[c]{0.4\linewidth}\centering
\begin{tikzpicture}[x=1mm, y=1mm]
  \draw[bkGray, line width=0.6pt] (0,0) circle (16);
  \fill[bkInk] (0,0) circle (0.9);
  \fill[bkInk] (-20:16) circle (0.9) (40:16) circle (0.9);
  \node[font=\sffamily\footnotesize, anchor=north west] at (-20:16.5) {P};
  \node[font=\sffamily\footnotesize, anchor=south west] at (40:16.8) {Q};
  \draw[kivk vec, bkBlue, line width=1.2pt] (-20:16) -- ++(70:12);
  \draw[kivk vec, bkOrange, line width=1.2pt] (40:16) -- ++(130:12);
  \draw[bkGray, line width=0.4pt, dashed] (0,0) -- (10:16);
  % 差（始点をそろえて、中心の近くに描く）
  \begin{scope}[shift={(-20,-32)}]
  \draw[kivk vec, bkBlue, line width=1.1pt] (0,0) -- (70:12);
  \draw[kivk vec, bkOrange, line width=1.1pt] (0,0) -- (130:12);
  \draw[kivk vec, bkRed, line width=1.1pt] ($(70:12)+(190:0.6)$) -- ($(130:12)+(10:0.6)$);
  \node[kivk lab, text=bkRed!80!black, anchor=south] at ($(100:12.6)$) {$\Delta\vec{v}$};
  \end{scope}
\end{tikzpicture}
\end{minipage}
\end{motto}

\bksubsection{6}{平行移動}

\label{bk:4-10:heikou}3-2では、点を「右へ3・上へ2」のように動かしました。この「動かし方」は、まさにベクトルです。では、点を1つだけでなく、図形全体を同じベクトルで動かしてみましょう。

\noindent\begin{minipage}[c]{0.54\linewidth}\centering
\begin{tikzpicture}[x=5.5mm, y=5.5mm]
  \draw[bkLine, line width=0.3pt] (0,0) grid (11,6);
  \filldraw[fill=bkBlue!15, draw=bkBlue!80!black, line width=0.8pt] (1,1) -- (4,1) -- (2,4) -- cycle;
  \filldraw[fill=bkOrange!18, draw=bkOrange!80!black, line width=0.8pt] (6,3) -- (9,3) -- (7,6) -- cycle;
  \draw[kivk vec, bkRed, line width=1pt] (1.15,1.06) -- (5.85,2.94);
  \draw[kivk vec, bkRed, line width=1pt] (4.15,1.06) -- (8.85,2.94);
  \draw[kivk vec, bkRed, line width=1pt] (2.15,4.06) -- (6.85,5.94);
  \node[font=\sffamily\footnotesize, anchor=north] at (1,0.9) {A};
  \node[font=\sffamily\footnotesize, anchor=north] at (4,0.9) {B};
  \node[font=\sffamily\footnotesize, anchor=east] at (1.9,4.1) {C};
  \node[font=\sffamily\footnotesize, anchor=north west] at (6.1,2.9) {A$'$};
  \node[font=\sffamily\footnotesize, anchor=north west] at (9,3) {B$'$};
  \node[font=\sffamily\footnotesize, anchor=south east] at (7,6) {C$'$};
\end{tikzpicture}
\end{minipage}\hfill
\begin{minipage}[c]{0.43\linewidth}\raggedright 
三角形ABCの3つの頂点を、どれも「右へ5・上へ2」のベクトル $\vec{v}$ で動かします。動かした先をA$'$、B$'$、C$'$とします（A$'$ は「エー・ダッシュ」と読みます）。\par\vspace{2pt}
赤い3本の矢印は、どれも $\vec{v}$ です。
\[ \overrightarrow{\mathrm{AA'}}=\overrightarrow{\mathrm{BB'}}=\overrightarrow{\mathrm{CC'}}=\vec{v} \]
\end{minipage}

\par\vspace{4pt}
三角形A$'$B$'$C$'$は、三角形ABCと形も大きさも同じ（合同、4-1）で、向きも変わっていません。4-1の「ずらす」動かし方です。図形のすべての点を同じベクトルで動かすことを、\textbf{平行移動}といいます。\bkidx{へいこういどう}{平行移動}

もう1つの見方をしてみましょう。駅の動く歩道に、荷物を3つ置きます。動く歩道が\qty{20}{m}進むと、3つの荷物はどれも、同じ向きに\qty{20}{m}動きます。荷物どうしの並び方は変わりません。これが平行移動です。

\bksubsection{7}{動いている人から見た速度：相対速度}

\label{bk:4-10:soutai}\bkivkstep{1}{電車の窓から}
電車に乗って、窓の外を見ています。となりの線路を、別の電車が走っています。

\noindent\begin{minipage}[c]{0.5\linewidth}\centering
\begin{tikzpicture}[x=1mm, y=1mm]
  \foreach \y in {0,12} {\draw[bkGray, line width=0.6pt] (0,\y-1) -- (56,\y-1);}
  \filldraw[fill=bkBlue!25, draw=bkBlue!80!black, line width=0.6pt, rounded corners=1mm] (2,12) rectangle (30,18);
  \filldraw[fill=bkOrange!30, draw=bkOrange!80!black, line width=0.6pt, rounded corners=1mm] (14,0) rectangle (42,6);
  \foreach \x in {5,11,17,23} {\fill[white] (\x,14.5) rectangle ++(4,2.5);}
  \foreach \x in {17,23,29,35} {\fill[white] (\x,2.5) rectangle ++(4,2.5);}
  \draw[kivk vec, bkBlue] (32,15) -- (46,15);
  \draw[kivk vec, bkOrange] (44,3) -- (61.5,3);
  \node[kivk lab, text=bkBlue!80!black, anchor=south] at (39,16.5) {東へ20};
  \node[kivk lab, text=bkOrange!80!black, anchor=south] at (52.5,4.5) {東へ25};
  \node[font=\sffamily\footnotesize, anchor=south] at (16,18.8) {A（自分）};
  \node[font=\sffamily\footnotesize, anchor=north] at (28,-1.8) {B};
\end{tikzpicture}
\end{minipage}\hfill
\begin{minipage}[c]{0.47\linewidth}\raggedright 
自分の電車Aは東へ\qty{20}{m/s}、となりの電車Bは東へ\qty{25}{m/s}で走っています（速度は地面に対するもの）。\par\vspace{2pt}
Aの窓から見ると、Bはどう見えるかな？
\end{minipage}

\par\vspace{4pt}
1秒ごとに、2つの電車の先頭の位置を調べます。0秒のとき、どちらも同じ所にいたとします。

\begin{bkcenter}
\begin{bktable}{colspec={Q[l,m,wd=40mm] *{4}{Q[c,m,wd=14mm]}}, rowsep=1.5pt}
時間（s） & 0 & 1 & 2 & 3 \\
Aの位置（m） & 0 & 20 & 40 & 60 \\
Bの位置（m） & 0 & 25 & 50 & 75 \\
BはAより何m前か & 0 & 5 & 10 & 15 \\
\end{bktable}
\bktabcaption{BはAより、1秒に\qty{5}{m}ずつ前へ出ていく。}
\end{bkcenter}

BはAより、1秒に\qty{5}{m}ずつ前へ出ていきます。だから、Aに乗っている人には、Bが東へ\qty{5}{m/s}でゆっくり追いこしていくように見えます。$25-20=5$ で、Bの速度からAの速度をひいた値です。

Bが反対に、西へ\qty{25}{m/s}で走っていたらどうでしょう。東を＋とすると、Bの速度は $-25$ なので、
\bkivkline{$(-25)-(+20)=-45$}
Aから見ると、Bは西へ\qty{45}{m/s}で、あっというまにすれちがっていきます。

このように、動いている人から見た相手の速度を、\textbf{相対速度}（そうたいそくど）といいます。\bkidx{そうたいそくど}{相対速度}
\bkivkline{\fbox{Aから見たBの相対速度 ＝ Bの速度 － Aの速度}}
\bkivkline{$\vec{v}_{\text{AB}}=\vec{v}_{\text{B}}-\vec{v}_{\text{A}}$}
自分（A）の速度をひくのは、「自分が止まって見えるように、まわりの全部から自分の速度を取りのぞく」ためです。自分の速度から自分の速度をひけば、$\vec{0}$ になります。

\bkivkstep{2}{平面の上の相対速度：車の窓から見た雨}
風のない日、雨つぶが真下へ\qty{6}{m/s}で落ちています。そこへ車が東へ\qty{8}{m/s}で走ってきます。車の窓から見ると、雨はどう見えるでしょう。

\noindent\begin{minipage}[c]{0.52\linewidth}\centering
\begin{tikzpicture}[x=1mm, y=1mm]
  % 雨の速度（下へ6）と、車の速度の反対（西へ8）
  \draw[kivk vec, bkBlue] (0,0) -- (0,-24);
  \draw[kivk vec, bkOrange] (0,-24) -- (-31.4,-24);
  \draw[kivk vec, bkRed] (-0.4,-0.6) -- (-31.4,-23.6);
  \fill[bkInk] (0,0) circle (1);
  \node[kivk lab, text=bkBlue!80!black, anchor=west, align=left] at (1.5,-12) {雨の速度\\下へ6};
  \node[kivk lab, text=bkOrange!80!black, anchor=north] at (-16,-25.5) {$-$（車の速度）：西へ8};
  \node[kivk lab, text=bkRed!80!black, anchor=south east, align=right] at (-17,-10) {車から見た\\雨の速度};
  \draw[bkInk, line width=0.5pt] (0,-6) arc (-90:-143.13:6);
  \node[font=\sffamily\footnotesize] at (-116.6:9.6) {$\theta$};
\end{tikzpicture}
\end{minipage}\hfill
\begin{minipage}[c]{0.45\linewidth}\raggedright 
車から見た雨の速度は、雨の速度から車の速度をひいたものです。\par\vspace{2pt}
「下へ6」に、車の速度の反対「西へ8」をたします。\par\vspace{2pt}
2辺が6と8の直角三角形の斜辺なので、大きさは $\sqrt{6^2+8^2}=10$ で、\qty{10}{m/s}です。
\end{minipage}

\par\vspace{4pt}
車の中から見ると、雨は\qty{10}{m/s}で、車の後ろの向きへななめに落ちてくるように見えます。たてからの角 $\theta$ は、$\tan\theta=\dfrac{8}{6}$ なので、電卓で約53°です。止まっているときはまっすぐ落ちていた雨のすじが、走り出すと横の窓にななめにつくのは、このためです。フロントガラスに雨が多く当たるのも、車から見ると雨が前からやってくるからです。

\bkfukidashi{mina}{自分の速度をひくと、自分から見た世界になるんだね。電車の窓の景色が後ろへ流れていくのも同じだ！}

\begin{taisetsu}
\def\tsl#1{\makebox[4.2\zw][l]{\sffamily\bfseries\small\color{bkDeep}#1}}
\tsl{ベクトル}向きと大きさをもつ量。矢印で表す。大きさは $|\vec{a}|$\\[3pt]
\tsl{和}矢印をつなぐ（平行四辺形の対角線）。$\overrightarrow{\mathrm{AB}}+\overrightarrow{\mathrm{BC}}=\overrightarrow{\mathrm{AC}}$\\[3pt]
\tsl{実数倍}$k\vec{a}$：$k>0$ なら同じ向き、$k<0$ なら反対向き。長さは $|k|$ 倍\\[3pt]
\tsl{差}$\vec{a}-\vec{b}=\vec{a}+(-\vec{b})$。$\overrightarrow{\mathrm{OA}}-\overrightarrow{\mathrm{OB}}=\overrightarrow{\mathrm{BA}}$（$\vec{b}$ の先から $\vec{a}$ の先へ）\\[3pt]
\tsl{速度}$\Delta\vec{v}=\vec{v}_{\text{あと}}-\vec{v}_{\text{まえ}}$\qquad 相対速度 $\vec{v}_{\text{AB}}=\vec{v}_{\text{B}}-\vec{v}_{\text{A}}$
\end{taisetsu}

\begin{kotoba}
\small
\bkword{ベクトル・スカラー}{向きと大きさをもつ量・大きさだけで決まる量。}
\bkword{始点・終点}{矢印のもと・先。$\overrightarrow{\mathrm{AB}}$ は始点A、終点B。}
\bkword{合力}{2つ以上の力を合わせて1つにした力。力のベクトルの和。}
\bkword{ゼロベクトル $\vec{0}$}{大きさ0のベクトル。向きは決まらない。}
\bkword{反対のベクトル}{大きさが同じで向きが反対のベクトル。$-\vec{a}$。}
\bkword{平行移動}{図形のすべての点を、同じベクトルで動かすこと。}
\bkword{相対速度}{動いている人から見た、相手の速度。相手の速度－自分の速度。}
\end{kotoba}

\begin{column}{「北の風」は、どっちへふく？}
天気予報の「北の風、風速\qty{5}{m/s}」は、\textbf{北からふいてくる}風のことです。空気は北から南へ流れます。天気図の風の記号も、風のふいてくる向きを表します。ところが、風をベクトル（空気の速度）として矢印でかくと、矢印は\textbf{南}を向きます。「北の風」の「北」は、矢印の向きではなく、矢印の始まる向きなのです。ベクトルで計算するときは、天気予報のことばを、空気の動く向きに読みかえてから使いましょう。風向（ふうこう）と風速の2つで1つの風を表すのは、まさに「向きと大きさ」です。
\end{column}

%% ---------------- 例題 ----------------
\bksection{reidai}{例題でたしかめよう}

\begin{reidai}[ベクトルとスカラー]
次の量は、ベクトルですか、スカラーですか。\par
（1）教室の気温\qty{22}{\degreeCelsius}\quad（2）北へ\qty{2}{km}の移動\quad（3）時速\qty{60}{km}という速さ\quad（4）机を右へ押す\qty{15}{N}の力\par
\kangaekata 「どっちへ」を言わないと決まらない量がベクトル、数と単位だけで言いつくせる量がスカラー。\par
\shiki（1）温度に向きはない。\quad（2）「北へ」と向きがある。\par
\hspace*{2.6\zw}（3）速さは、向きを言っていない。\quad（4）「右へ」と向きがある。\par
\kotae （1）スカラー\quad（2）ベクトル\quad（3）スカラー\quad（4）ベクトル
\end{reidai}
（3）は、「北へ時速\qty{60}{km}」のように向きをそえれば、速度（ベクトル）になります。

\begin{jibun}
次の量は、ベクトルですか、スカラーですか。\par
（1）ペットボトルの水の体積\qty{500}{mL}\quad（2）東へ\qty{3}{m/s}の速度\quad（3）質量\qty{40}{kg}\quad（4）物を真上へ持ち上げる\qty{20}{N}の力
\bkjans{（1）\nobreak スカラー（2）\nobreak ベクトル（3）\nobreak スカラー（4）\nobreak ベクトル}
\end{jibun}

\begin{reidai}[直角に引く2つの力の合力]
荷物を、東へ\qty{12}{N}、北へ\qty{5}{N}の力で同時に引きます。合力の大きさと、東から北へ何度回った向きかを求めましょう（電卓を使ってよい）。\par
\kangaekata 2つの力を2辺とする長方形の対角線が合力。大きさは三平方の定理（4-4）、角は $\tan$（4-6）。\par
\shiki 大きさは $\sqrt{12^2+5^2}=\sqrt{144+25}=\sqrt{169}=13$\par
\hspace*{2.6\zw}東からの角を $\theta$ とすると、$\tan\theta=\dfrac{5}{12}=0.4166\cdots$\par
\hspace*{2.6\zw}電卓（DEG）で $\tan^{-1}(5\div12)=22.61\cdots$\par
\kotae \qty{13}{N}、東から約23°北へ回った向き
\end{reidai}

\begin{jibun}
荷物を、東へ\qty{15}{N}、北へ\qty{8}{N}の力で同時に引きます。合力の大きさと、東から北へ何度回った向きかを求めましょう。
\bkjans{$\sqrt{15^2+8^2}=\sqrt{289}=\qty{17}{N}$。$\tan^{-1}\dfrac{8}{15}\approx28.1°$ で、東から約28°北へ回った向き}
\end{jibun}

\bkivkmidashi{電卓で角を求めるときは、DEG を確かめる}
\noindent\begin{minipage}[c]{0.56\linewidth}\raggedright
電卓の $\tan^{-1}$ は、ふつう SHIFT（または 2nd）のキーを押してから tan のキーを押して使います。角が度で出るように、画面に DEG（度）の印が出ていることを確かめましょう。RAD（ラジアン、4-2）のままだと、同じ計算で0.3947…と出てきます。これは $0.3947\times\dfrac{180}{\pi}\approx22.6$ で、22.6°をラジアンで表した数です。
\end{minipage}\hfill
\begin{minipage}[c]{0.4\linewidth}\centering
\begin{tikzpicture}[x=1mm, y=1mm]
  \bkdisplay{(0,14)}{46}{22.61986495}
  \node[font=\sffamily\footnotesize\bfseries, anchor=west, text=bkInk] at (2,22.5) {DEG};
  \bkdisplay{(0,0)}{46}{0.394791119}
  \node[font=\sffamily\footnotesize\bfseries, anchor=west, text=bkInk] at (2,8.5) {RAD};
\end{tikzpicture}
\end{minipage}

\begin{reidai}[角をなす2つの力の合力]
\qty{3}{N}の力と\qty{5}{N}の力が、同じ点に、なす角60°ではたらいています。合力の大きさを求めましょう。\par
\kangaekata 平行四辺形をかき、片方の力を平行に動かしてつなぐと、2辺3と5、そのあいだの角が $180°-60°=120°$ の三角形ができる（数学の道具③の段5）。余弦定理（4-7）で向かいの辺を求める。\par
\shiki $(\text{合力})^2=3^2+5^2-2\times3\times5\times\cos120°$\par
\hspace*{4.6\zw}$=9+25-30\times(-0.5)=9+25+15=49$\par
\hspace*{2.6\zw}合力 $=\sqrt{49}=7$\par
\kotae \qty{7}{N}
\end{reidai}

\begin{jibun}
\qty{7}{N}の力と\qty{8}{N}の力が、同じ点に、なす角60°ではたらいています。合力の大きさを求めましょう。
\bkjans{$7^2+8^2-2\times7\times8\times\cos120°=49+64+56=169$。$\sqrt{169}=\qty{13}{N}$}
\end{jibun}

\bkivkmidashi{作図でも確かめよう}
\noindent\begin{minipage}[c]{0.5\linewidth}\raggedright
例題3を、ノートに作図して確かめます。\qty{1}{N}を\qty{1}{cm}として、\qty{5}{cm}と\qty{3}{cm}の矢印を、分度器で60°開いてかきます。平行四辺形を作って対角線を定規ではかると、約\qty{7}{cm}になります。計算の答え\qty{7}{N}と合います。作図は答えの見当をつけるのに、計算は正確な値を出すのに役立ちます。
\end{minipage}\hfill
\begin{minipage}[c]{0.46\linewidth}\centering
\begin{tikzpicture}[x=0.85mm, y=0.85mm]
  \draw[bkGray, line width=0.5pt, dashed] (15,25.98) -- (65,25.98) -- (50,0);
  \draw[kivk vec, bkBlue, line width=1.2pt] (0,0) -- (50,0);
  \draw[kivk vec, bkOrange, line width=1.2pt] (0,0) -- (15,25.98);
  \draw[kivk vec, bkRed, line width=1.2pt] (0.9,0.36) -- (64.1,25.62);
  \node[kivk lab, text=bkBlue!80!black, anchor=north] at (25,-1.5) {\qty{5}{cm}};
  \node[kivk lab, text=bkOrange!80!black, anchor=east] at (6,14) {\qty{3}{cm}};
  \node[kivk lab, text=bkRed!80!black, anchor=north west] at (38,13) {約\qty{7}{cm}};
\end{tikzpicture}
\end{minipage}

\begin{reidai}[方眼の上の実数倍・和・差]
\noindent\begin{minipage}[c]{0.6\linewidth}\raggedright 
右の方眼の $\vec{a}$、$\vec{b}$ について、次のベクトルを「右へ○・上へ○」の形で答えましょう。（1）は大きさも求めましょう。左・下へ進むときは「左へ」「下へ」と書きます。\par
（1）$\vec{a}+\vec{b}$\quad（2）$2\vec{a}$\quad（3）$-\vec{b}$\quad（4）$\vec{a}-\vec{b}$
\end{minipage}\hfill
\begin{minipage}[c]{0.36\linewidth}\centering
\begin{tikzpicture}[x=5mm, y=5mm]
  \draw[bkLine, line width=0.3pt] (0,0) grid (8,4);
  \draw[kivk vec, bkBlue] (0.5,0.5) -- (2.5,1.5);
  \draw[kivk vec, bkOrange] (4.5,0.5) -- (5.5,3.5);
  \node[kivk lab, text=bkBlue!80!black, anchor=south east] at (1.3,1.35) {$\vec{a}$};
  \node[kivk lab, text=bkOrange!80!black, anchor=east] at (4.7,2.2) {$\vec{b}$};
\end{tikzpicture}
\end{minipage}\par\vspace{2pt}
\kangaekata $\vec{a}$ は「右へ2・上へ1」、$\vec{b}$ は「右へ1・上へ3」。つなぐと、横の動きどうし、たての動きどうしがたされる（3-2）。\par
\shiki（1）$\vec{a}$ に $\vec{b}$ をつなぐ：右へ $2+1=3$、上へ $1+3=4$\par
\hspace*{4.6\zw}大きさは $\sqrt{3^2+4^2}=5$\par
\hspace*{2.6\zw}（2）$\vec{a}$ を2本つなぐ：右へ $2\times2=4$、上へ $1\times2=2$\par
\hspace*{2.6\zw}（3）$\vec{b}$ と反対向き：左へ1、下へ3\par
\hspace*{2.6\zw}（4）$\vec{a}+(-\vec{b})$：横は $2-1=1$、たては $1-3=-2$\par
\hspace*{4.6\zw}右へ1、下へ2\par
\kotae （1）右へ3・上へ4、大きさ5\quad（2）右へ4・上へ2\quad（3）左へ1・下へ3\quad（4）右へ1・下へ2
\end{reidai}
（4）を図でかくと、$\vec{a}$ と $\vec{b}$ の始点をそろえたとき、$\vec{b}$ の先から $\vec{a}$ の先へ向かう矢印になっています。ノートに方眼をかいて、確かめてみましょう。

\begin{jibun}
$\vec{a}$ が「右へ3・上へ1」、$\vec{b}$ が「右へ1・上へ2」のとき、次のベクトルを「右へ○・上へ○」の形で答えましょう。（1）は大きさも求めましょう。\par
（1）$\vec{a}+\vec{b}$\quad（2）$2\vec{a}$\quad（3）$-\vec{b}$\quad（4）$\vec{a}-\vec{b}$
\bkjans{（1）\nobreak 右へ4・上へ3、大きさ5（2）\nobreak 右へ6・上へ2（3）\nobreak 左へ1・下へ2（4）\nobreak 右へ2・下へ1}
\end{jibun}

\begin{reidai}[一直線上の速度の変化]
東へ\qty{8}{m/s}で転がってきたボールが、壁に当たって、西へ\qty{6}{m/s}ではね返りました。速度の変化 $\Delta\vec{v}$ を求めましょう。\par
\kangaekata 東を＋と決めて、向きを符号で表す（2-2）。速度の変化は「あと－まえ」。\par
\shiki まえの速度 $+8$、あとの速度 $-6$\par
\hspace*{2.6\zw}$\Delta v=(-6)-(+8)=-14$\par
\hspace*{2.6\zw}負なので西向き\par
\kotae 西へ\qty{14}{m/s}
\end{reidai}

\begin{jibun}
南へ\qty{20}{m/s}で飛んできたテニスボールを、ラケットで北へ\qty{30}{m/s}で打ち返しました。速度の変化を求めましょう。
\bkjans{北を＋として $(+30)-(-20)=+50$。北へ\qty{50}{m/s}}
\end{jibun}

\begin{reidai}[向きが変わったときの速度の変化]
自転車が、速さ\qty{6}{m/s}のまま、北向きから東向きへ曲がりました。速度の変化 $\Delta\vec{v}$ の大きさと向きを求めましょう。\par
\kangaekata $\Delta\vec{v}=\vec{v}_{\text{あと}}-\vec{v}_{\text{まえ}}=\vec{v}_{\text{あと}}+(-\vec{v}_{\text{まえ}})$。始点をそろえてかく（図\ref{bk:4-10:fig-dv}）。\par
\shiki $\vec{v}_{\text{あと}}$ は東へ6、$-\vec{v}_{\text{まえ}}$ は南へ6\par
\hspace*{2.6\zw}東へ6と南へ6の和なので、向きは南東\par
\hspace*{2.6\zw}大きさは $\sqrt{6^2+6^2}=\sqrt{72}=6\sqrt2\approx8.49$\par
\kotae 南東の向きに約\qty{8.5}{m/s}
\end{reidai}
速さは\qty{6}{m/s}のままなのに、速度の変化は0ではありません。

\begin{jibun}
車が、速さ\qty{15}{m/s}のまま、南向きから東向きへ曲がりました。速度の変化 $\Delta\vec{v}$ の大きさと向きを求めましょう。
\bkjans{東へ15と北へ15の和。北東の向きに $15\sqrt2\approx\qty{21.2}{m/s}$}
\end{jibun}

\begin{reidai}[電車から見た電車の速度]
電車Aが東へ\qty{15}{m/s}で走っています。（1）電車Bが東へ\qty{25}{m/s}で走っているとき、（2）電車Bが西へ\qty{25}{m/s}で走っているとき、Aから見たBの相対速度を求めましょう。\par
\kangaekata Aから見たBの相対速度＝Bの速度－Aの速度。東を＋とする。\par
\shiki（1）$(+25)-(+15)=+10$\par
\hspace*{2.6\zw}（2）$(-25)-(+15)=-40$\par
\kotae （1）東へ\qty{10}{m/s}\quad（2）西へ\qty{40}{m/s}
\end{reidai}

\begin{jibun}
自転車Aが北へ\qty{5}{m/s}で走っています。（1）自転車Bが北へ\qty{3}{m/s}で走っているとき、（2）自転車Bが南へ\qty{4}{m/s}で走っているとき、Aから見たBの相対速度を求めましょう。
\bkjans{（1）\nobreak $3-5=-2$ で南へ\qty{2}{m/s}（Aから見ると、Bは後ろへ下がっていく）（2）\nobreak $(-4)-5=-9$ で南へ\qty{9}{m/s}}
\end{jibun}

\bkfukidashi{sora}{同じ向きに走っていても、自分より遅い相手は、後ろへ下がっていくように見えるんだね。}

\begin{reidai}[車から見た雨の速度]
雨つぶが真下へ\qty{5}{m/s}で落ちています。人が自転車で東へ\qty{5}{m/s}で走ると、人から見た雨の速度の大きさと向きはどうなるでしょう。\par
\kangaekata 人から見た雨の速度＝雨の速度－人の速度＝雨の速度＋（人の速度の反対）。\par
\shiki 「下へ5」と「西へ5」の和\par
\hspace*{2.6\zw}大きさは $\sqrt{5^2+5^2}=5\sqrt2\approx7.07$\par
\hspace*{2.6\zw}2辺が等しい直角三角形なので、たてから45°西へかたむく\par
\kotae 大きさ約\qty{7.1}{m/s}。たてから45°だけ西（人の後ろの向き）へかたむいて落ちる
\end{reidai}

\begin{jibun}
雨つぶが真下へ\qty{6}{m/s}で落ちています。車が北へ\qty{8}{m/s}で走ると、車から見た雨の速度の大きさと、たてから何度かたむくかを求めましょう。
\bkjans{「下へ6」と「南へ8」の和。$\sqrt{6^2+8^2}=\qty{10}{m/s}$。$\tan^{-1}\dfrac{8}{6}\approx53°$ だけ南（車の後ろ）へかたむく}
\end{jibun}

\begin{reidai}[平行移動した点の座標]
三角形ABCの頂点は $\mathrm{A}(-2,\ 1)$、$\mathrm{B}(1,\ 0)$、$\mathrm{C}(0,\ 3)$ です。この三角形を、「右へ4・上へ3」のベクトル $\vec{v}$ で平行移動した三角形A$'$B$'$C$'$の頂点の座標を求めましょう。また、1つの頂点が動いた距離を求めましょう。\par
\kangaekata 平行移動では、どの点も同じベクトルで動く（数学の道具⑥）。右へ4なら $x$ 座標に4をたし、上へ3なら $y$ 座標に3をたす（3-2）。動いた距離は $|\vec{v}|$。\par
\shiki A$'$：$(-2+4,\ 1+3)=(2,\ 4)$\par
\hspace*{2.6\zw}B$'$：$(1+4,\ 0+3)=(5,\ 3)$\par
\hspace*{2.6\zw}C$'$：$(0+4,\ 3+3)=(4,\ 6)$\par
\hspace*{2.6\zw}動いた距離 $|\vec{v}|=\sqrt{4^2+3^2}=5$\par
\kotae A$'(2,\ 4)$、B$'(5,\ 3)$、C$'(4,\ 6)$。どの頂点も5だけ動く
\end{reidai}
3つの頂点がどれも同じ5だけ、同じ向きに動くので、三角形の形も大きさも向きも変わりません。ノートに方眼をかいて、2つの三角形をかいて確かめましょう。

\begin{jibun}
三角形ABCの頂点は $\mathrm{A}(1,\ 2)$、$\mathrm{B}(4,\ -1)$、$\mathrm{C}(-1,\ 0)$ です。この三角形を、「右へ3・下へ2」のベクトルで平行移動した三角形A$'$B$'$C$'$の頂点の座標と、1つの頂点が動いた距離を求めましょう。
\bkjans{A$'(4,\ 0)$、B$'(7,\ -3)$、C$'(2,\ -2)$。動いた距離 $\sqrt{3^2+2^2}=\sqrt{13}\approx3.61$}
\end{jibun}

%% ---------------- まちがえやすいところ ----------------
\bksection{tsumazuki}{まちがえやすいところをチェック}

\begin{tsumazuki}{大きさをそのままたしてしまう}
\bkng{東へ\qty{6}{N}と北へ\qty{8}{N}の合力は $6+8=\qty{14}{N}$}
\bkok{合力は長方形の対角線。$\sqrt{6^2+8^2}=\qty{10}{N}$}
大きさをそのままたしてよいのは、同じ向きのときだけです。向きがちがうと、引っぱり合う分だけ、合力は小さくなります。
\end{tsumazuki}

\begin{tsumazuki}{差の矢印の向きを逆にかく}
\bkng{$\vec{a}-\vec{b}$ を、$\vec{a}$ の先から $\vec{b}$ の先へかく}
\bkok{$\vec{a}-\vec{b}$ は、$\vec{b}$ の先から $\vec{a}$ の先へ}
$\vec{b}$ に $\vec{a}-\vec{b}$ をたすと $\vec{a}$ になるはずです。$\vec{b}$ の先から出発して $\vec{a}$ の先に着く矢印でなければ、たして $\vec{a}$ になりません。
\end{tsumazuki}

\begin{tsumazuki}{速さが同じなら、速度は変わっていないと思う}
\bkng{速さ\qty{10}{m/s}のまま曲がったので、$\Delta\vec{v}=\vec{0}$}
\bkok{向きが変わったので、$\Delta\vec{v}\neq\vec{0}$（90°曲がると $10\sqrt2$）}
速度はベクトルなので、大きさ（速さ）が同じでも、向きが変われば別の速度です。
\end{tsumazuki}

\begin{tsumazuki}{ベクトルとその大きさを、同じ記号で書く}
\bkng{$\vec{F}=\qty{40}{N}$}
\bkok{$|\vec{F}|=\qty{40}{N}$（向きもいうなら「$\vec{F}$ は東へ\qty{40}{N}」）}
$\vec{F}$ は向きをもつ矢印、$|\vec{F}|$ はただの数です。矢印と数を「＝」で結ぶことはできません。
\end{tsumazuki}

\begin{tsumazuki}{相対速度で、ひく順番を逆にする}
\bkng{東へ\qty{20}{m/s}の電車Aから見た、東へ\qty{25}{m/s}の電車Bの速度は $20-25=-5$ で西へ\qty{5}{m/s}}
\bkok{Aから見たBは、Bの速度－Aの速度。$25-20=5$ で東へ\qty{5}{m/s}}
ひくのは「見ている自分（A）」の速度です。Aの窓から見ると、速いBは前へ出ていくので、東向きでなければおかしいと、答えを見て確かめましょう。
\end{tsumazuki}

\begin{column}{「ベクトル」という名前}
「ベクトル」という名前を数学で使いはじめたのは、アイルランドのハミルトンで、180年ほど前のことです。その後、アメリカのギブズとイギリスのヘビサイドが、電気と磁気の研究（マクスウェルの電磁気学）を計算しやすくするために、いま私たちが使っている矢印のたし算・ひき算や、4-12・4-14で学ぶ2種類のかけ算を整理しました。力の平行四辺形はずっと昔から知られていましたが、「向きと大きさをもつ量」をまとめて1つの文字で計算できるようになったのは、意外に新しいのです。
\end{column}

%% ---------------- 物理・工学 ----------------
\begin{butsuri}
\begin{itemize}[leftmargin=1.2\zw, itemsep=2pt, topsep=0pt, label={\color{bkPurple}●}]
\item \textbf{クレーンの2本のワイヤ}\par
\noindent\begin{minipage}[c]{0.62\linewidth}
クレーンで重い荷物をつるとき、フックから2本のワイヤを荷物の両はしにかけます。2本のワイヤの力の合力が、荷物の重さとつり合います。2本のあいだの角（つり角）を広げると、合力を同じに保つために、1本1本の力は大きくなります。つり角120°では、1本のワイヤに、荷物の重さと同じ大きさの力がかかります（2本の力と重さの3つの矢印で、正三角形ができるから）。現場では、つり角を原則として60°以内にするように決められています。
\end{minipage}\hfill
\begin{minipage}[c]{0.34\linewidth}\centering
\begin{tikzpicture}[x=1mm, y=1mm]
  \draw[bkInk, line width=0.8pt] (0,26) -- (0,20);
  \draw[bkInk, line width=1pt] (0,20) arc (90:-200:1.6);
  \draw[bkGray!80!black, line width=0.9pt] (0,18.4) -- (-13,2) (0,18.4) -- (13,2);
  \filldraw[fill=bkOrange!30, draw=bkOrange!70!black, line width=0.6pt] (-16,-8) rectangle (16,2);
  \draw[bkInk, line width=0.5pt] ($(0,18.4)+(-128.4:6)$) arc (-128.4:-51.6:6);
  \node[font=\sffamily\footnotesize] at (0,9.5) {つり角};
  \draw[kivk vec, bkPurple, line width=1.1pt] (0,-3) -- (0,-15);
  \node[font=\sffamily\footnotesize, anchor=west] at (1,-12) {重さ};
\end{tikzpicture}
\end{minipage}
\par\vspace{2pt}
{\centering\small
\begin{tblr}{colspec={Q[l,m] *{4}{Q[c,m,wd=12mm]}}, hlines={0.5pt,bkLine}, vlines={0.5pt,bkLine}, row{1}={bg=bkTheme!12, font=\sffamily\small}, rowsep=1pt, colsep=4pt}
つり角 & 0° & 60° & 90° & 120° \\
1本の力÷荷物の重さ & 0.50 & 0.58 & 0.71 & 1.00 \\
\end{tblr}\par}
\item \textbf{横風の中を飛ぶ飛行機}\par
飛行機は空気の中を進みます。地面から見た飛行機の速度は、「空気に対する飛行機の速度」と「風の速度」の和です。横風がふくと、機首の向きと実際に進む向きがずれるので、パイロットは機首を風上へ少し向けて、まっすぐ滑走路へ進みます（4-11で計算します）。
\item \textbf{ロボットの手の移動}\par
工場のロボットは、手先を「いまの位置から、目標の位置まで」動かします。動かす向きと距離は、目標の位置から、いまの位置をひいたベクトル（あと－まえ）です。
\end{itemize}
\end{butsuri}

\begin{column}{斜面の首かざり}
\noindent\begin{minipage}[c]{0.66\linewidth}
400年あまり前、オランダのステヴィンは、右のような図を考えました。三角の台に、同じ玉をつないだ首かざりをかけます。もし首かざりがひとりでに回り出したら、いつまでも回り続けることになりますが、そんなことは起こりません。下にたれた部分は左右対称なので、それを取りのぞいても、斜面の上の玉はつり合っているはずです。ステヴィンはここから、斜面の上の物を支える力のきまりを見つけ、力を三角形で合わせる考えへ進みました。
\end{minipage}\hfill
\begin{minipage}[c]{0.3\linewidth}\centering
\begin{tikzpicture}[x=1mm, y=1mm]
  \fill[bkGray!30] (0,0) -- (36,0) -- (24,18) -- cycle;
  \draw[bkInk!70, line width=0.6pt] (0,0) -- (36,0) -- (24,18) -- cycle;
  \foreach \t in {0.08,0.22,...,0.96} {\fill[bkBlue!70] ($(0,0)!\t!(24,18)+(-1.1,1.5)$) circle (1.2);}
  \foreach \t in {0.15,0.45,0.75} {\fill[bkBlue!70] ($(36,0)!\t!(24,18)+(1.2,0.8)$) circle (1.2);}
  \draw[bkBlue!70, line width=0.6pt] (-1.5,1.5) .. controls (6,-14) and (30,-14) .. (37.5,0.8);
  \foreach \t in {0.12,0.3,0.5,0.7,0.88} {\fill[bkBlue!70] ($(-1.5,1.5)!\t!(37.5,0.8)+(0,{-10*sin(180*\t)})$) circle (1.2);}
\end{tikzpicture}
\end{minipage}
\end{column}

%% ---------------- 章末問題 ----------------
\bkendproblems
\bkproblemgroup{基本}{2}

\bkproblem 次の量を、ベクトルとスカラーに分けましょう（\bkreidairef{1}）。\\
（1）北東へ\qty{100}{m}の移動\quad（2）\qty{30}{s}という時間\quad（3）下向きの\qty{9.8}{N}の重力\quad（4）\qty{2}{L}という体積
\bkpans{ベクトル：（1）（3）。スカラー：（2）（4）}

\bkproblem 荷物を、東へ\qty{9}{N}、北へ\qty{12}{N}の力で同時に引きます。合力の大きさと、東から北へ何度回った向きかを求めましょう（\bkreidairef{2}）。
\bkpans{$\sqrt{9^2+12^2}=\sqrt{225}=\qty{15}{N}$。$\tan^{-1}\dfrac{12}{9}\approx53.1°$ で、東から約53°北へ回った向き}

\bkproblem \qty{5}{N}の力と\qty{3}{N}の力が、なす角60°ではたらいています。合力の大きさを求めましょう（\bkreidairef{3}）。
\bkpans{$5^2+3^2-2\times5\times3\times\cos120°=49$。\qty{7}{N}}

\bkproblem $\vec{a}$ が「右へ1・上へ2」、$\vec{b}$ が「右へ3・下へ1」のとき、次のベクトルを「右へ○・上へ○」の形で答えましょう（\bkreidairef{4}）。\quad（1）$\vec{a}+\vec{b}$\quad（2）$3\vec{a}$\quad（3）$\vec{a}-\vec{b}$
\bkpans{（1）\nobreak 右へ4・上へ1（2）\nobreak 右へ3・上へ6（3）\nobreak 左へ2・上へ3}

\bkproblem 東へ\qty{12}{m/s}で飛んできたボールを、バットで西へ\qty{30}{m/s}で打ち返しました。速度の変化を求めましょう（\bkreidairef{5}）。
\bkpans{東を＋として $(-30)-(+12)=-42$。西へ\qty{42}{m/s}}

\bkproblem 電車Aが東へ\qty{20}{m/s}で走っています。電車Bが（1）東へ\qty{12}{m/s}、（2）西へ\qty{12}{m/s}で走るとき、Aから見たBの相対速度を求めましょう（\bkreidairef{7}）。
\bkpans{（1）\nobreak $12-20=-8$ で西へ\qty{8}{m/s}（2）\nobreak $(-12)-20=-32$ で西へ\qty{32}{m/s}}

\bkproblem 荷物を、東へ\qty{20}{N}、北へ\qty{21}{N}の力で同時に引きます。合力の大きさと、東から北へ何度回った向きかを求めましょう（\bkreidairef{2}）。
\bkpans{$\sqrt{20^2+21^2}=\sqrt{841}=\qty{29}{N}$。$\tan^{-1}\dfrac{21}{20}\approx46.4°$ で、東から約46°北へ回った向き}

\bkproblem 三角形ABCの頂点は $\mathrm{A}(0,\ 0)$、$\mathrm{B}(3,\ 1)$、$\mathrm{C}(1,\ 4)$ です。「左へ2・上へ5」のベクトルで平行移動した三角形A$'$B$'$C$'$の頂点の座標と、1つの頂点が動いた距離を求めましょう（\bkreidairef{9}）。
\bkpans{A$'(-2,\ 5)$、B$'(1,\ 6)$、C$'(-1,\ 9)$。動いた距離 $\sqrt{2^2+5^2}=\sqrt{29}\approx5.39$}

\bkproblemgroup{チャレンジ}{3}

\bkproblem 速さ\qty{20}{m/s}のまま、西向きから北向きへ曲がった車の、速度の変化の大きさと向きを求めましょう（\bkreidairef{6}）。
\bkpans{北へ20と東へ20の和。北東の向きに $20\sqrt2\approx\qty{28.3}{m/s}$}

\bkproblem \qty{6}{N}の力と\qty{10}{N}の力が、なす角60°ではたらいています。合力の大きさを求めましょう（\bkreidairef{3}）。
\bkpans{$36+100+60=196$ で\qty{14}{N}}

\bkproblem 雨つぶが真下へ\qty{4}{m/s}で落ちています。人が自転車で西へ\qty{3}{m/s}で走ると、人から見た雨の速度の大きさと、たてから何度かたむくかを求めましょう（\bkreidairef{8}）。
\bkpans{「下へ4」と「東へ3」の和。$\sqrt{4^2+3^2}=\qty{5}{m/s}$。$\tan^{-1}\dfrac{3}{4}\approx37°$ だけ東（人の後ろ）へかたむく}

\par\newpage
\bkproblem 雨つぶが真下へ\qty{8}{m/s}で落ちています。車が東へ\qty{6}{m/s}で走ると、車から見た雨の速さと、たてからかたむく角を求めましょう（\bkreidairef{8}）。
\bkpans{「下へ8」と「西へ6」の和。$\sqrt{8^2+6^2}=\qty{10}{m/s}$。$\tan^{-1}\dfrac{6}{8}\approx37°$ だけ西（車の後ろ）へかたむく}

\bkproblem \qty{5}{N}の力と\qty{16}{N}の力が、なす角60°ではたらいています。合力の大きさを求めましょう（\bkreidairef{3}）。
\bkpans{$25+256+80=361$ で\qty{19}{N}}

\bkproblem 自転車が、速さ\qty{12}{m/s}のまま、北向きから西向きへ曲がりました。速度の変化の大きさと向きを求めましょう（\bkreidairef{6}）。
\bkpans{西へ12と南へ12の和。南西の向きに $12\sqrt2\approx\qty{17.0}{m/s}$}

\begin{furikaeri}
\begin{checklist}
\item 向きと大きさをもつ量を矢印と記号で表せる。\bkfurref{bk:4-10:kigou}
\item 矢印をつないで、合力の大きさを求められる。\bkfurref{bk:4-10:wa}
\item 差をかいて、速度の変化や相対速度を求められる。\bkfurref{bk:4-10:sa}
\end{checklist}
\end{furikaeri}

\bkfukidashi{delta}{矢印の計算は、作図だけだと大変だね。次の章では、矢印を2つの数の組にして、計算だけで答えを出すよ。}
